\documentclass[11pt,a4paper]{article}
\usepackage[T1]{fontenc}
\usepackage[utf8]{inputenc}
\usepackage{lmodern,amsmath,amssymb,amsthm,mathtools,mathrsfs}
\usepackage[margin=26mm,headheight=15pt]{geometry}
\usepackage{microtype,booktabs,longtable,tabularx,array,enumitem}
\usepackage{xcolor,listings,fancyhdr,titlesec,xurl}
\definecolor{inkblue}{RGB}{28,55,83}
\definecolor{quietgray}{RGB}{83,91,99}
\definecolor{papergray}{RGB}{247,248,250}
\usepackage[colorlinks=true,linkcolor=inkblue,citecolor=inkblue,urlcolor=inkblue,
 bookmarksnumbered=true,pdfencoding=auto]{hyperref}
\pagestyle{fancy}
\fancyhf{}
\fancyhead[L]{\small\color{quietgray}Ten External Problems: Proofs and Limits}
\fancyhead[R]{\small\color{quietgray}Keio Research Notes}
\fancyfoot[C]{\thepage}
\renewcommand{\headrulewidth}{0.3pt}
\titleformat{\section}{\Large\bfseries\color{inkblue}}{\thesection}{0.7em}{}
\titleformat{\subsection}{\large\bfseries\color{inkblue}}{\thesubsection}{0.65em}{}
\titleformat{\subsubsection}{\normalsize\bfseries}{\thesubsubsection}{0.6em}{}
\setlength{\emergencystretch}{3em}
\setlength{\parskip}{3pt}
\setlist{itemsep=3pt,topsep=5pt}
\setcounter{tocdepth}{1}
\numberwithin{equation}{section}
\newtheorem{theorem}{Theorem}[section]
\newtheorem{lemma}[theorem]{Lemma}
\newtheorem{proposition}[theorem]{Proposition}
\newtheorem{corollary}[theorem]{Corollary}
\theoremstyle{definition}
\newtheorem{definition}[theorem]{Definition}
\newtheorem{hypothesis}[theorem]{Hypothesis}
\theoremstyle{remark}
\newtheorem{remark}[theorem]{Remark}
\newtheorem{example}[theorem]{Example}
\newcommand{\A}{\overline{\mathbb Q}}
\newcommand{\Q}{\mathbb Q}
\newcommand{\Z}{\mathbb Z}
\newcommand{\N}{\mathbb N}
\newcommand{\R}{\mathbb R}
\newcommand{\C}{\mathbb C}
\newcommand{\Er}{\mathcal E}
\newcommand{\K}{\mathcal K}
\newcommand{\Ein}{\operatorname{Ein}}
\newcommand{\Eone}{E_1}
\newcommand{\Tr}{\operatorname{Tr}}
\newcommand{\Norm}{\operatorname{N}}
\newcommand{\vp}{v_p}
\newcommand{\floor}[1]{\left\lfloor #1\right\rfloor}
\newcommand{\ceil}[1]{\left\lceil #1\right\rceil}
\DeclareMathOperator{\Frac}{Frac}
\DeclareMathOperator{\trdeg}{trdeg}
\DeclareMathOperator{\den}{den}
\DeclareMathOperator{\ord}{ord}
\DeclareMathOperator{\Res}{Res}
\lstdefinestyle{certificate}{language=Python,basicstyle=\ttfamily\fontsize{8.5}{10.2}\selectfont,emptylines=0,
 keywordstyle=\color{inkblue}\bfseries,commentstyle=\color{quietgray},
 stringstyle=\color{inkblue},backgroundcolor=\color{papergray},
 frame=single,rulecolor=\color{quietgray!35},framerule=0.3pt,
 breaklines=true,breakatwhitespace=false,columns=fullflexible,
 keepspaces=true,showstringspaces=false,tabsize=4,
 aboveskip=9pt,belowskip=9pt,xleftmargin=5pt,xrightmargin=5pt}

\hypersetup{pdftitle={From Mills Clocks and Euler's Constant to Ten External Problems},
 pdfauthor={Keio and Shaitetsu},
 pdfsubject={Two full external answers, one analytic identity extension, and seven partial targets}}
\fancyhead[L]{\small\color{quietgray}Ten External Problems: Proofs and Limits}
\DeclareMathOperator{\Log}{Log}
\setcounter{tocdepth}{2}


\title{\color{inkblue}\textbf{From Mills Clocks and Euler's Constant\\to Ten External Problems}\\[7pt]
\large Complete Answers, Exact Identities,\\and Partial Arithmetic Theorems}
\author{Keio \and Shaitetsu}
\date{5 November 2026\\[3pt]\small A single research manuscript with full proofs and a status register}
\begin{document}
\maketitle
\thispagestyle{empty}
\begin{abstract}
We apply modular return, integer-tail, monodromy, and residue methods
from a preceding investigation of Euler's constant and shifted Mills
constants to ten externally formulated targets. Two questions of Saito
are answered with their full quantifiers by reconstructing earlier
corpus arguments: a non-reversible inhomogeneous recurrence produces
composite Pisot-floor windows, and every fixed integer shift of the
dyadically indexed Fibonacci sequence is composite on a positive-density
set of indices. We also resolve the analytic-identity component of a
recent Meijer-function extension problem: explicit differential operators
evaluate to $1$ and $\gamma+\Log z$ at every argument, with a higher
zeta-polynomial hierarchy and an exact telescoping kernel. Further
theorems classify $E$-value units in the field of all algebraic exponential
values with one $\operatorname{Ein}$ value, prove irrationality for
power-growth subfamilies of the monotone Erd\H{o}s 270 problem, classify
primitive harmonic-lattice seeds, and give an all-rank scalar approximation
to $\gamma+\log z$. Three additional prime-clock targets receive precise
partial results. Every external target is matched to its original source
and its actual outcome. Seven of the ten targets remain unresolved in
their full stated scope. Neither irrationality of $\gamma$ nor
unconditional transcendence of the classical Mills constant is proved.
\end{abstract}
\vfill
\noindent\textbf{Keywords.} Mills constants; Euler--Mascheroni constant;
Pisot numbers; Lucas sequences; $E$-function values; monodromy;
factorial series; Meijer $G$-functions; conservative matrix fields.

\medskip
\noindent\textbf{Reading convention.} A complete proof in this manuscript
is distinct from a claim of historical priority or external peer review.
The source register identifies prior corpus arguments, new derivations
in this investigation, and remaining hypotheses.
\clearpage
\tableofcontents
\clearpage

\section{The ten targets and their exact outcomes}
\label{ext:register}

This paper follows \cite{Integrated2026}, which organized ten
gamma-related and ten Mills-related items from the research dialogue.
The present task is different: select ten actual questions outside that
discussion and determine what the developed methods prove about them.
We keep each external question's quantifiers visible. An auxiliary
irrational series is not a proof that its base is irrational; a
classification on one field is not a classification of every $E$-value;
and scalar convergence is not generic matrix-state convergence.

Our notation $\xi((C_n))$ means the least $A>1$ for which
$\lfloor A^{C_n}\rfloor$ is prime for every $n\ge1$, when that
least value exists. In particular
\[
 M=\xi((3^n)),\qquad \xi_2=\xi((2^n)).
\]
The four shifted clocks $3^n\pm1$ and $3^n\pm n$ motivated the
preceding investigation. The full external answers below use their
modular-clock and arithmetic methods; they do not require an unproved
identity expressing $\gamma$ in terms of those four constants.

\subsection{Status and provenance}

The register uses three outcomes. \emph{Full answer} means all
quantifiers in the named external problem are proved here.
\emph{Identity component resolved} means the precisely selected
analytic component of a broader future-work paragraph is proved.
\emph{Partial; original target open here} means a complete auxiliary
theorem is provided, but the original external conclusion is not.
No statement that a source problem was globally unsolved immediately
before this manuscript is inferred merely from an unsuccessful search.
Sources and versions were checked on 5 November 2026.

The two full Saito answers already existed in corpus manuscripts
[007] and [008], dated 5 September 2026; we reconstruct their proofs
and strengthen the exposition rather than count them as discoveries
first made today. The conditional Mills-gap and three-adic-window
arguments likewise have identified corpus predecessors. The residue,
field-intersection, denominator-window, and scalar-concentration
derivations are developed in the present investigation; literature
priority for equivalent results is not asserted.

\small
\begin{longtable}{@{}p{0.055\textwidth}p{0.365\textwidth}p{0.51\textwidth}@{}}
\toprule
ID & External target and source & Outcome in this manuscript\\
\midrule\endhead
\hyperref[ext:gamma]{T1} & Irrationality of Euler's constant $\gamma$; current approximation context \cite{CMF2026}. & \textbf{Partial.} Exact primitive-seed classification, height cost, and a conditional radical-membership obstruction. The irrationality question is not closed.\\[7pt]
\hyperref[exm:mills]{T2} & Classical Mills transcendence, through Saito's cubic-Pisot alternative, Question 1.7 \cite{Saito2025}. & \textbf{Partial.} A hypothetical algebraic tower has a transcendental gap exponent and divisibility constraints. The cubic-Pisot alternative is not excluded.\\[7pt]
\hyperref[exm:binary]{T3} & Is the least binary-clock constant $\xi_2$ irrational? Question 1.6 \cite{Saito2025}. & \textbf{Partial.} The reciprocal floor series is irrational for every real base $A>1$. This does not decide the arithmetic of $\xi_2$.\\[7pt]
\hyperref[exm:cubic]{T4} & Infinitely many composite floors on $3^n$ for every cubic Pisot number; Problem 1.1 \cite{SaitoIntervals}. & \textbf{Partial.} Nonintegral three-adic trace limits give simultaneous composite windows. Exceptional integral limits remain.\\[7pt]
\hyperref[exm:nonreversible]{T5} & Find a non-reversible ILRS working for every Pisot number; Problem 1.7 \cite{SaitoIntervals}. & \textbf{Full answer, reconstructed from [008].} $R(n)=ba^n\pm n$ works, with simultaneous windows and an extension to Salem numbers.\\[7pt]
\hyperref[exm:fibonacci]{T6} & Infinitely many composite $F_{2^n}+h$ for every integer $h$; Problem 1.8 \cite{SaitoIntervals}. & \textbf{Full answer, reconstructed from [007].} A positive lower density of composite indices, with a general Lucas-sequence version.\\[7pt]
\hyperref[ev:units-section]{T7} & Are all units of the $E$-value ring $\alpha e^\beta$? Conjecture 1.2 \cite{FR2026}. & \textbf{Partial.} Complete classification over the field of all algebraic exponential values and one nonzero algebraic Ein argument; the global conjecture remains.\\[7pt]
\hyperref[es:270]{T8} & Irrationality of the factorial series for every nondecreasing $f\to\infty$; monotone Erd\H{o}s 270 variant \cite{CrKo2025}. & \textbf{Partial.} All $f(n)=\lfloor cn^\alpha\rfloor$, $c>0$, $\alpha>3/2$, and a boundary range; all integer-valued polynomials of degree at least two.\\[7pt]
\hyperref[cmf:extension]{T9} & Extend the Meijer identity to $\gamma+\log z$; analytic component of Section 9.2(1) \cite{CMF2026}. & \textbf{Identity component resolved.} Every $p\ge2$ and principal-branch $z$, a zeta-polynomial hierarchy, and the exact kernel of weighted telescoping.\\[7pt]
\hyperref[cmf:convergence]{T10} & Generic dual convergence for all admissible ranks and trajectories; Conjecture 1 \cite{CMF2026}. & \textbf{Partial.} A separate all-rank scalar convergence theorem and integer dilations of the known matrix family. General dual-state isolation is not proved.\\
\bottomrule
\end{longtable}
\normalsize

Problem and theorem numbers for Saito refer to the arXiv versions
specified in the bibliography; journal versions need not use the same
numbering. The unproved Riemann Hypothesis is not inserted as an
assumption into an allegedly unconditional conclusion. The implication
from the Density Hypothesis to Mills transcendence is an existing
external theorem \cite{SaitoVariants}.

\subsection{The common proof mechanisms}

There are three useful transfers. On prime clocks, an accessible residue
class turns a single fixed divisor into infinitely many composite
values. In rational series, a common denominator turns a nonzero small
tail into an impossible integer. For $E$-values, analytic continuation
produces algebraic embeddings whose incompatible ramification and
denominator patterns force descent. The Meijer calculation uses a
different exact mechanism: an infinite residue sum telescopes to one
boundary residue. Each method proves a finite, checkable implication;
the remaining challenge is to establish the hypothesis needed by the
original target.


\section{Euler's constant: exact arithmetic of adaptive seeds}
\label{ext:gamma}

The first target is the irrationality of
\[
 \gamma=\lim_{n\to\infty}(H_n-\log n),\qquad
 H_n=\sum_{j=1}^n\frac1j.
\]
The current recurrence construction of \cite{CMF2026} does not settle
this question. This section gives an exact arithmetic classification
of the two-dimensional image lattice appearing in the preceding
integrated manuscript \cite{Integrated2026}. Its purpose is to test
whether selecting a new seed at every depth creates an irrationality
proof. The classification shows exactly what this selection costs.

\subsection{Primitive points on a harmonic shear lattice}

For integers $D\ge1$ and $C$ with $\gcd(C,D)=1$, define
\[
 A_{D,C}=\begin{pmatrix}D&C\\0&D\end{pmatrix},\qquad
 \Lambda_{D,C}=A_{D,C}\mathbb Z^2.
\]
The normalized image in the rank-three trajectory examined in
\cite[Section 8]{Integrated2026} has this form with
$D=\operatorname{den}(3H_n)$ and $C=3DH_n$. The theorem below is
independent of that identification: it holds for every displayed
integral shear, and its proof needs no analytic estimate for a matrix
recurrence.

\begin{theorem}[Exact primitive-seed classification]
\label{ext:seed-classification}
Let $(p,q)\in\mathbb Z^2$ satisfy $q>0$ and $\gcd(p,q)=1$. Put
\[
 h=\gcd(Dp-Cq,Dq).
\]
Then $h\mid D^2$, and the unique primitive integer input whose image
is a positive multiple of $(p,q)$ is
\begin{equation}\label{ext:primitive-seed}
 (a,b)=\left(\frac{Dp-Cq}{h},\frac{Dq}{h}\right).
\end{equation}
Its image and the gcd of the two image coordinates are
\begin{equation}\label{ext:seed-output}
 A_{D,C}(a,b)^T=\frac{D^2}{h}(p,q)^T,
 \qquad \gcd(Da+Cb,Db)=\frac{D^2}{h}.
\end{equation}
If $\gcd(q,D)=1$, then $h=1$ and this output gcd is exactly $D^2$.
\end{theorem}
\begin{proof}
The defining gcd makes $(a,b)$ integral and primitive. Also $h$
divides
\[
 D(Dp-Cq)+C(Dq)=D^2p,
 \qquad D(Dq)=D^2q.
\]
Bezout's identity for $p,q$ therefore gives $h\mid D^2$.
Matrix multiplication proves \eqref{ext:seed-output}; primitiveness
of $(p,q)$ identifies its gcd. Any input on the same image ray lies
on the rational line spanned by $(Dp-Cq,Dq)$. The only primitive
integer point on its positive ray is the one in
\eqref{ext:primitive-seed}.

Finally suppose $\gcd(q,D)=1$ and a prime $\ell$ divides $h$.
If $\ell\mid q$, then $\ell\nmid Dp$, contradicting
$\ell\mid Dp-Cq$. Otherwise $\ell\mid D$, since $\ell\mid Dq$,
and then $\ell\nmid Cq$, giving the same contradiction.
Thus $h=1$.
\end{proof}

Every rational slope is consequently represented. This is an exact
surjectivity statement about slopes, not an estimate on the quality
of approximation to $\gamma$.

\begin{corollary}[The cost of a chosen rational approximation]
\label{ext:seed-cost}
Suppose $D_n=\operatorname{den}(3H_n)$, $C_n=3D_nH_n$, and
$p_n/q_n$ is a bounded sequence of reduced rational numbers with
$q_n>0$. With $h_n$ as in Theorem~\ref{ext:seed-classification},
the effective primitive input satisfies
\[
 \max(|a_n|,|b_n|)
 \sim \frac{3D_nq_n\log n}{h_n}.
\]
The image linear form and its primitive normalization are exactly
\[
 (D_nb_n)\gamma-(D_na_n+C_nb_n)
   =\frac{D_n^2}{h_n}(q_n\gamma-p_n),
\]
and $q_n\gamma-p_n$, respectively.
\end{corollary}
\begin{proof}
Equation \eqref{ext:primitive-seed} gives
$a_n=(D_nq_n/h_n)(p_n/q_n-3H_n)$ and $b_n=D_nq_n/h_n$.
Since $H_n/\log n\to1$, the first coordinate eventually dominates
and has the claimed size. The linear-form identity follows by
substitution in \eqref{ext:seed-output}.
\end{proof}

Here ``effective input'' means the two free coordinates of the image
lattice. The norm of a three-coordinate preimage under the original
CMF matrix is a further question. A bound for a fixed matrix seed is
not automatically uniform for the depth-dependent inputs just
constructed.

\subsection{What an irrationality transfer must still prove}

\begin{proposition}[Primitive integer-form criterion]
\label{ext:gamma-criterion}
If integers $p_j,q_j$ satisfy
\[
 0<|q_j\gamma-p_j|\longrightarrow0,
\]
then $\gamma$ is irrational. Obtaining only
$p_j/q_j\to\gamma$ does not supply this conclusion.
\end{proposition}
\begin{proof}
If $\gamma=u/v$ with $v>0$, each nonzero $q_j\gamma-p_j$ has
absolute value at least $1/v$. For the final distinction, even a
rational number $x$ has rational approximations $p_j/q_j\to x$;
division by $q_j$ can hide the nonvanishing size of its integer
linear form.
\end{proof}

The seed classification neither proves the strict nonvanishing nor
the decay required by this proposition. It shows that arbitrary
adaptive seed selection simply re-encodes the selected rational
approximation, with its exact arithmetic multiplier. The decisive
remaining estimate must concern the primitive form after this
multiplier has been removed.

\subsection{A second, algebraic transfer obstruction}

Let $\mathcal E$ be the ring of $E$-values, $\mathcal K$ its fraction
field, and put
\[
 \eta=\operatorname{Ein}(1),\qquad
 \delta=\int_0^\infty\frac{e^{-t}}{1+t}\,dt.
\]
The usual exponential-integral identity \cite[Section 6.2]{NISTDLMF}
gives
\begin{equation}\label{ext:gamma-ein}
 \eta=\gamma+\delta/e.
\end{equation}
The classification in Theorem~\ref{ev:one-Ein-field}, proved below,
turns this relation into the following exact conditional obstruction.

\begin{corollary}[A radical membership obstruction]
\label{ext:gamma-radical}
If $\gamma$ is algebraic, then
\[
 \delta\in\mathcal E\setminus\mathcal E^\times,
 \qquad \delta^{-1}\notin\mathcal E,
 \qquad \delta^{1/m}\notin\mathcal K\quad(m\ge2).
\]
Here the last assertion holds for every choice of an $m$th root.
Consequently, proving $\delta^{1/m}\in\mathcal K$ for even one
$m\ge2$ would prove that $\gamma$ is transcendental.
\end{corollary}
\begin{proof}
Write $A=\overline{\mathbb Q}$, $\mathcal B=A[e^\beta:\beta\in A]$,
$\mathcal L=\operatorname{Frac}(\mathcal B)$, and
$F=\mathcal L(\eta)$. Under the displayed hypothesis,
$\delta=e(\eta-\gamma)$ is a nonconstant polynomial in $\eta$
over $\mathcal B$. Theorem~\ref{ev:one-Ein-field} makes it a
non-unit of $\mathcal E$. The same theorem says that $F$ is
relatively algebraically closed in $\mathcal K$. An $m$th root
of $\delta$ in $\mathcal K$ would therefore lie in $F$.
In the rational function field $\mathcal L(\eta)$, the valuation
at the prime $\eta-\gamma$ of $\delta$ equals one. The valuation
of an $m$th power is divisible by $m$, a contradiction.
\end{proof}

No membership statement for such a radical is established here.
The corollary specifies a possible transfer theorem; it is not a
proof that its antecedent can be met. In particular, the irrationality
of $\gamma$ remains unresolved in this manuscript.


% Self-contained section fragment: no preamble, labels prefixed exm:.
% Requires amsmath, amssymb, amsthm, hyperref and standard theorem environments.
% Keys: Saito2025, SaitoVariants, SaitoIntervals, Schneider1934,
% VMRC007, VMRC008, VMRC010, VMRC042.

\section{Five external questions about prime-representing constants}
\label{exm:section}

The five targets below retain the quantifiers of their published
formulations. Two are answered completely by arguments already present
in the supplied research corpus, which are independently reconstructed
here. The other three remain unresolved; the accompanying theorems are
partial results or consequences of additional hypotheses.

\subsection{Question 1.6: the least square-clock constant}
\label{exm:binary}

For an integer \(c\ge2\), let
\[
 \xi_c=\min\{A>1:\lfloor A^{c^n}\rfloor
                  \text{ is prime for every }n\ge1\}.
\]
Existence is known. Saito's Question~1.6 in
\cite{Saito2025} asks whether \(\xi_2\) is rational or irrational.
The question is not answered here. The integer-tail mechanism does,
however, give a complete result at the critical growth rate \(2^n\).
This auxiliary result is proved directly; no claim of historical
priority is made.

\begin{theorem}[Irrationality of the reciprocal square-clock series]
\label{exm:square-series}
For every real number \(A>1\), the convergent series
\[
 S(A)=\sum_{n=1}^{\infty}\frac1{\lfloor A^{2^n}\rfloor}
\]
is irrational.
\end{theorem}

\begin{proof}
Write \(a_n=\lfloor A^{2^n}\rfloor\) and
\(P_n=\prod_{j=1}^n a_j\). The initial values may equal \(1\);
all values are positive and \(a_n\to\infty\). Since
\[
 0\le1-\frac{a_n}{A^{2^n}}<A^{-2^n},
\]
the infinite product
\(\Omega=\prod_{n\ge1}a_n/A^{2^n}\)
converges to a strictly positive number. Moreover,
\(a_{n+1}\ge a_n^2\). The first term dominates each tail:
\[
 \sum_{j>n}\frac1{a_j}
 =\frac1{a_{n+1}}\bigl(1+o(1)\bigr),
\]
because the remaining terms are at most
\(\sum_{\ell\ge1}a_{n+1}^{-2^\ell}
 =O(a_{n+1}^{-2})\).
As \(\sum_{j=1}^n2^j=2^{n+1}-2\), it follows that
\begin{equation}\label{exm:square-limit}
 P_n\sum_{j>n}\frac1{a_j}\longrightarrow A^{-2}\Omega>0.
\end{equation}
Suppose \(S(A)=u/v\), where \(u,v\) are positive integers.
Then
\[
 I_n=vP_n\sum_{j>n}\frac1{a_j}
 =uP_n-v\sum_{j=1}^n\frac{P_n}{a_j}
\]
is a positive integer. By \eqref{exm:square-limit}, this integer
sequence converges, and hence is eventually an integer \(m>0\).
The exact identity
\(I_{n+1}=a_{n+1}I_n-vP_n\)
gives \(m(a_{n+1}-1)=vP_n\). Comparing successive indices yields
\[
 a_{n+2}-1=a_{n+1}(a_{n+1}-1).
\]
Thus \(a_{n+2}=a_{n+1}^2-a_{n+1}+1<a_{n+1}^2\) eventually,
contradicting \(a_{n+2}\ge a_{n+1}^2\).
\end{proof}

The theorem applies to \(A=\xi_2\), but also to \(A=2\).
Consequently its conclusion cannot determine the arithmetic
nature of \(A\). A proof of Question~1.6 still needs an implication
from a rational hypothetical base to additional restrictions on its
prime-floor sequence. The irrationality of \(S(\xi_2)\) supplies no
such implication by itself.

\subsection{Question 1.7: the cubic Pisot alternative for Mills' constant}
\label{exm:mills}

Write \(M=\xi_3\). Saito's Question~1.7 in \cite{Saito2025}
asks whether \(M^{3^m}\) can be a cubic Pisot number for some
positive integer \(m\). A negative answer would prove \(M\)
transcendental. The later reduction in
\cite[Theorem~1.7]{SaitoVariants} shows that, if \(M\) is algebraic,
some such power \(\beta=M^{3^m}\) has three real conjugates
\begin{equation}\label{exm:saito-reduction}
 \beta,\ -\rho,\ \tau,\qquad
 0<\sigma:=|\tau|<\rho<1,\qquad
 \rho\le\min(\sigma^{17/23},\beta^{-17/40}).
\end{equation}
This is an external theorem. The conditional gap theorem below
is a corpus result from \cite{VMRC010}, included with its proof because
it specifies what an unconditional transcendence argument must exclude.

\begin{theorem}[Arithmetic of a hypothetical algebraic Mills gap]
\label{exm:gap}
Assume \(M\) algebraic and choose \(\beta\) as in
\eqref{exm:saito-reduction}. Set
\[
 p_n=\lfloor\beta^{3^n}\rfloor,\qquad
 g_n=p_{n+1}-p_n^3,\qquad N=|\operatorname{Norm}(\beta)|.
\]
On a sufficiently late tail,
\[
 g_n\mid g_{n+1},\qquad
 \frac{9g_{n+1}}{g_n^3}\longrightarrow1,\qquad
 \frac{g_n}{p_n}\longrightarrow\infty.
\]
There is a transcendental real number
\(\kappa=(2+\log\rho/\log\beta)/3\) such that
\begin{equation}\label{exm:gap-asymptotic}
 g_n\sim3(p_n^3)^\kappa,\qquad
 \frac12+\frac{\log N}{6\log\beta}
 <\kappa<\frac{21}{40}.
\end{equation}
\end{theorem}

\begin{proof}
The sum of the two small conjugate powers is
\(-\rho^{3^n}+\tau^{3^n}\in(-1,0)\) eventually. The integer
trace is therefore \(p_n=\operatorname{Tr}(\beta^{3^n})\).
At a fixed index write \(X,Y,Z\) for the three conjugate powers
and put \(p=X+Y+Z\), \(b=XY+XZ+YZ\), \(c=XYZ\).
These are integers. Cubing the roots gives
\[
 p'=p^3-3bp+3c,\quad
 b'=b^3-3cbp+3c^2,\quad c'=c^3.
\]
In particular, \(h_n=g_n/3=c-bp\) and
\[
 h_n=-(X+Y)(X+Z)(Y+Z),\qquad
 h_{n+1}=h_n
 \prod_{\{U,V\}\subset\{X,Y,Z\}}(U^2-UV+V^2).
\]
The multiplier is a symmetric polynomial with integer coefficients
in the algebraic integers \(X,Y,Z\), so is an integer.
It is positive because the three numbers are nonzero and real.
This proves divisibility.

The inequalities between the conjugate moduli give
\[
 p_n\sim\beta^{3^n},\qquad
 h_n\sim(\beta^2\rho)^{3^n}.
\]
They imply the gap asymptotic and
\(9g_{n+1}/g_n^3\to1\). Also
\(1\le N=\beta\rho\sigma<\beta\rho^2\).
Thus \(\rho>\sqrt{N/\beta}\), which gives the strict lower bound
for \(\kappa\); it also gives \(\beta\rho>1\), proving
\(g_n/p_n\to\infty\). The last inequality in
\eqref{exm:saito-reduction} gives \(\kappa\le21/40\).

Let \(t=\log\rho/\log\beta\). If \(t\) were rational, its
negativity would imply \(\beta^u\rho^v=1\) for positive
integers \(u,v\). A transitive cubic Galois group contains a
three-cycle. Choose its orientation to send
\(\beta\mapsto-\rho\mapsto\tau\). Applying it to the last
identity and taking absolute values gives
\(\rho^u\sigma^v=1\), impossible. Hence \(t\) is irrational.
If \(t\) were algebraic, the Gelfond--Schneider
theorem would make the positive real value
\(\beta^t=\rho\) transcendental \cite{Schneider1934}.
Thus \(t\), and therefore \(\kappa\), is transcendental.
The rational upper endpoint \(21/40\) cannot occur.
\end{proof}

\begin{corollary}[An exact conditional route to transcendence]
\label{exm:gap-criterion}
Let \(q_n=\lfloor M^{3^n}\rfloor\) be the actual Mills primes.
If the limit
\[
 \lim_{n\to\infty}
 \frac{\log(q_{n+1}-q_n^3)}{\log(q_n^3)}
\]
exists and is an algebraic real number, then \(M\) is
transcendental.
\end{corollary}
\begin{proof}
Under algebraicity, deleting finitely many indices identifies
the limit with the transcendental \(\kappa\) in
Theorem~\ref{exm:gap}, a contradiction.
\end{proof}

No argument here proves that the limit in the corollary is
algebraic or excludes the alternative in
\eqref{exm:saito-reduction}. The unconditional Mills question
remains open in this paper. The implication
RH \(\Rightarrow M\) transcendental, and the stronger implication
from the Density Hypothesis, are already Saito's
Theorem~1.8 in \cite{SaitoVariants}; they are not a new resolution.

\subsection{Problem 1.1: composite values on the pure cubic clock}
\label{exm:cubic}

Problem~1.1 of \cite{SaitoIntervals} concerns every cubic Pisot
number \(\beta\), including a nonreal conjugate pair:
must \(\lfloor\beta^{3^n}\rfloor\) be composite infinitely often?
Here cubic clock refers to the exponent \(3^n\), not to a
pure cubic number field. We establish a substantial conditional
family but do not answer the universal question. This
simultaneous-window argument is an audited corpus result
\cite{VMRC042}.

\begin{lemma}[A three-adic trace limit]\label{exm:three-limit}
Let \(f(X)=X^3-uX^2+vX-N\) be a monic integral cubic.
Let \(T_n\) be the \(3^n\)-th power sum of its roots,
let \(B_n\) be the second elementary symmetric function of
those powered roots, and put \(E_n=N^{3^n}\).
There is \(\lambda\in\mathbb Z_3\) such that
\[
 \lambda-T_n\in3^{n+1}\mathbb Z_3\qquad(n\ge0).
\]
The limit depends only on \((u,v,N)\pmod3\).
\end{lemma}
\begin{proof}
Newton's identities give the integral polynomial iteration
\((T_{n+1},B_{n+1},E_{n+1})=F(T_n,B_n,E_n)\), where
\[
 F(x,y,z)=(x^3-3xy+3z,\ y^3-3zxy+3z^2,\ z^3).
\]
It is the identity on \(\mathbb F_3^3\). If two
integer vectors agree modulo \(3^r\), \(r\ge1\), their
images agree modulo \(3^{r+1}\). Indeed, every first
derivative has coefficients divisible by \(3\), and terms
of order at least two have the required divisibility.
Successive iterates differ by a vector in \(3^{n+1}\mathbb Z^3\).
Completeness gives the limit and its precision. The same
contraction applied to two initial states congruent modulo
\(3\) proves the final assertion.
\end{proof}

\begin{theorem}[A common progression of composite trace shifts]
\label{exm:windows}
Suppose \(T_n\to+\infty\) in the real topology. Let
\(\mathcal H\subset\mathbb Z\) be finite. If
\(\lambda+h\notin\{-1,1\}\) for every \(h\in\mathcal H\),
then some arithmetic progression of indices contains only
composite values \(T_n+h\), simultaneously for all
\(h\in\mathcal H\), after finitely many initial exceptions.
\end{theorem}
\begin{proof}
Choose a common large \(n_0\ge1\) for which
\(Q_h=T_{n_0}+h>1\) for every shift. When
\(3\nmid\lambda+h\), also require
\[
 v_3(Q_h-1)<n_0,\qquad v_3(Q_h+1)<n_0.
\]
This is possible because \(\lambda+h-1\) and
\(\lambda+h+1\) are nonzero, and
Lemma~\ref{exm:three-limit} eventually fixes their valuations.
If \(3\mid\lambda+h\), then \(3\mid T_n+h\) at every index.

For each other shift, some prime factor \(\ell_h\) of \(Q_h\)
must satisfy
\begin{equation}\label{exm:three-budget}
 v_3(\ell_h^2-1)\le n_0-1.
\end{equation}
Otherwise every prime factor would be \(1\) or \(-1\)
modulo \(3^{n_0}\): for a prime different from \(3\),
at most one of its predecessor and successor is divisible
by \(3\). Their product would then make \(Q_h\) equal
to \(1\) or \(-1\) modulo \(3^{n_0}\), a contradiction.
Here \(\ell_h\ne3\), since \(3\nmid Q_h\).

Every nonzero root of \(f\) over \(\mathbb F_{\ell_h}\)
belongs to an extension of degree at most three; its order
divides
\[
 L_h=\operatorname{lcm}(\ell_h-1,\ell_h^2-1,\ell_h^3-1).
\]
For \(\ell_h\equiv1\pmod3\), the last factor has
valuation \(v_3(\ell_h-1)+1\); for
\(\ell_h\equiv-1\pmod3\), the only nonzero valuation among
the three factors is \(v_3(\ell_h+1)\).
Thus \eqref{exm:three-budget} implies \(v_3(L_h)\le n_0\).
After raising the roots to \(3^{n_0}\), their nonzero
orders are prime to \(3\). Choose a period \(D_h\ge1\)
for cubing on these finitely many cyclic groups. Then
\[
 T_{n_0+jD_h}\equiv T_{n_0}\pmod{\ell_h}\qquad(j\ge0).
\]
The power-sum argument also covers zero and repeated roots,
with multiplicities.
Take a common multiple \(D\) of all \(D_h\), or \(D=1\)
if no such primes were needed. At \(n=n_0+jD\), each
trace shift has its fixed prime divisor. Real growth makes
every divisor proper eventually.
\end{proof}

\begin{corollary}[Pisot floor windows, including nonreal conjugates]
\label{exm:floor-windows}
Let \(\beta\) be any cubic Pisot number. If its trace limit
\(\lambda\) is not an ordinary integer, then for every
\(H\ge0\), all the numbers
\[
 \lfloor\beta^{3^n}\rfloor+h,\qquad |h|\le H,
\]
are simultaneously composite on an eventual arithmetic
progression of indices. Their common index set has positive
lower natural density.
\end{corollary}
\begin{proof}
The two small conjugate powers have real sum tending to zero.
The floor is therefore eventually \(T_n\) or \(T_n-1\),
whether or not that sum changes sign. Apply
Theorem~\ref{exm:windows} to every integer shift in
\([-H-1,H]\). Nonintegrality ensures
\(\lambda+h\ne\pm1\) for all these shifts.
\end{proof}

For a concrete coefficient class, irreducibility of the cubic
modulo \(3\) guarantees the required nonintegrality.
In the unramified degree-three extension of \(\mathbb Q_3\),
the roots are units. A unit is its Teichm\"uller root of unity
times a principal unit. The principal-unit part tends to \(1\)
under \(3^n\)-th powering. The three Teichm\"uller parts are
permuted by Frobenius, so their sum is \(\lambda\).
Irreducibility forces their orders to be \(13\) or \(26\).
For a primitive thirteenth root \(\zeta\),
\[
 \lambda=\pm(\zeta^a+\zeta^{3a}+\zeta^{9a}),
 \qquad a\not\equiv0\pmod{13}.
\]
This cannot be rational: otherwise a polynomial of degree
at most \(12\), with at most four nonzero coefficients,
would vanish at \(\zeta\), but it cannot be a nonzero
multiple of \(1+X+\cdots+X^{12}\).
For example \(X^3-X-1\) has no root modulo \(3\).
Its real root \(\beta>1\) is Pisot, because the nonreal
conjugate pair has modulus \(\beta^{-1/2}<1\).

The remaining obstruction is explicit. In a hypothetical algebraic
Mills tower, the floor is eventually \(T_n\). Its primality
at every index forces \(\lambda\in\{-1,1\}\) by
Theorem~\ref{exm:windows}. These two possibilities are not
excluded, so the partial result does not close Problem~1.1.

\subsection{Problem 1.7: a non-reversible recurrence for every Pisot number}
\label{exm:nonreversible}

Problem~1.7 in \cite{SaitoIntervals} asks for a non-reversible
inhomogeneous integer linear recurrence \(R(n)\), chosen
independently of the Pisot number \(\alpha\), for which
\(\lfloor\alpha^{R(n)}\rfloor\) is composite infinitely often
for every \(\alpha\). The following answers the full formulation.
It reconstructs \cite{VMRC008}, rather than presenting a
discovery made in this paper. An elementary trace-period proof
avoids any theorem on distribution of prime values.

\begin{lemma}[Every residue occurs on a progression]\label{exm:residues}
Let \(a\ge2\), \(b\ge1\) be integers and
\(\varepsilon\in\{-1,1\}\). For every \(q\ge1\) and \(c\in\mathbb Z\),
the solutions to
\[
 R(n):=ba^n+\varepsilon n\equiv c\pmod q
\]
include all sufficiently large indices in an arithmetic
progression.
\end{lemma}
\begin{proof}
For \(q>1\), factor \(q=uv\), where \(u\) contains exactly
the prime-power factors supported on primes dividing \(a\),
and \(\gcd(a,v)=1\). The sequence \(a^n\bmod q\) is
eventually periodic with a period \(t<q\): take
\(t=\operatorname{ord}_v(a)\le\varphi(v)<v\) if \(v>1\),
and \(t=1\) otherwise.

We prove arbitrarily late solutions exist by strong induction
on \(q\). The case \(q=1\) is immediate. Put
\(g=\gcd(t,q)<q\). By induction choose an arbitrarily
large \(n\), beyond the preperiod, with \(R(n)\equiv c\pmod g\).
Since \(\gcd(\varepsilon t,q)=g\), there is a nonnegative
\(j\) with \(\varepsilon jt\equiv c-R(n)\pmod q\).
Periodicity gives \(R(n+jt)\equiv c\pmod q\).
Finally \(R(n)\bmod q\) is eventually periodic with period
\(\operatorname{lcm}(q,t)\), so any sufficiently late solution
generates the claimed progression.
\end{proof}

\begin{theorem}[Full solution of the non-reversible recurrence question]
\label{exm:nonrev-theorem}
For every \(a\ge2,b\ge1,\varepsilon\in\{-1,1\}\), the
sequence \(R(n)=ba^n+\varepsilon n\), indexed by \(n\ge1\), is a positive,
increasing, non-reversible inhomogeneous integer linear recurrence.
No alternative integer recurrence representation makes it
reversible. For every Pisot or Salem number \(\alpha\) and
every \(H\ge0\), all
\[
 \lfloor\alpha^{R(n)}\rfloor+h,\qquad |h|\le H,
\]
are simultaneously composite on an eventual arithmetic progression
of indices. In particular \(R(n)=3^n+n\) answers
Saito's Problem~1.7.
\end{theorem}
\begin{proof}
The recurrence is
\[
 R(n+2)=(a+1)R(n+1)-aR(n)+\varepsilon(1-a).
\]
Positivity and strict increase follow from
\(ba^n-n\ge2^n-n\ge1\) and
\(R(n+1)-R(n)=b(a-1)a^n+\varepsilon\ge1\).
To exclude a hidden reversible representation, suppose
\(Q(E)R\) is constant, where \(E\) is the shift and \(Q\)
a monic integer polynomial. The coefficient of \(a^n\)
forces \(Q(a)=0\), and that of \(n\) forces \(Q(1)=0\).
Hence \((X-a)(X-1)\) divides \(Q\) in \(\mathbb Z[X]\).
Its constant coefficient is divisible by \(a\), so the
trailing recurrence coefficient \(-Q(0)\) cannot equal
\(1\) or \(-1\).

Let \(d\) be the degree of \(\alpha\), and set
\(T_m=\operatorname{Tr}(\alpha^m)\). This is an integer
recurrence sequence. All conjugates except \(\alpha\) have
modulus at most \(1\), so
\[
 T_m\longrightarrow+\infty,\qquad
 |\lfloor\alpha^m\rfloor-T_m|\le d.
\]
The recurrence matrix has nonzero integer determinant,
equal up to sign to the constant coefficient of the minimal
polynomial. Modulo primes not dividing this determinant,
the state is purely periodic. There are finitely many other
primes, and each finite state space gives eventual periodicity.
Thus an index \(s_0\), independent of the prime, exists
after which \(T_m\) is periodic modulo every prime, with
a prime-dependent period.

Put \(J=H+d\), and choose \(m_0\ge s_0\) so large that
\(T_{m_0}+j>1\) for all integers \(|j|\le J\).
For each \(j\), choose a prime divisor \(\ell_j\) of this
integer. Let \(L\) be a common multiple of the finitely
many trace periods modulo \(\ell_j\).
For \(m\ge m_0\) with \(m\equiv m_0\pmod L\),
\[
 \ell_j\mid T_m+j\qquad(|j|\le J).
\]
Real growth makes these divisors proper eventually.
Lemma~\ref{exm:residues} supplies an eventual progression
of \(n\) with \(R(n)\equiv m_0\pmod L\), and
\(R(n)\to\infty\). Finally each
\(\lfloor\alpha^{R(n)}\rfloor+h\), \(|h|\le H\),
equals \(T_{R(n)}+j\) for some integer \(|j|\le H+d\).
The whole required window is composite on the same progression.
\end{proof}

The distinction from Problem~1.1 is exact: the added term
\(\varepsilon n\) makes every residue accessible and is used
in the induction. The proof gives no such surjectivity
for \(R(n)=3^n\). Dropping the added term drops an essential
hypothesis.

\subsection{Problem 1.8: every shifted dyadic Fibonacci sequence}
\label{exm:fibonacci}

Problem~1.8 of \cite{SaitoIntervals} asks whether
\(F_{2^r}+h\) is composite infinitely often for every fixed
integer \(h\), with the usual Fibonacci initial values
\(F_0=0,F_1=1\). The result below answers the complete
question, with positive lower density. It independently
reconstructs the existing corpus proof \cite{VMRC007}.

More generally, fix a positive odd integer \(P\), and let
\[
 U_0=0,\quad U_1=1,\quad U_{n+2}=PU_{n+1}+U_n,\qquad
 D=P^2+4.
\]
For \(h\in\mathbb Z\), define the finite integer
\begin{equation}\label{exm:threshold}
 B_{P,h}=\max\left\{
 v_2\!\left(D(h-1)^2+3\right),
 v_2\!\left(D(h+1)^2+3\right)\right\}.
\end{equation}

\begin{theorem}[Prime return and composite Fibonacci shifts]
\label{exm:fib-theorem}
If \(r\ge B_{P,h}+1\) and \(p=U_{2^r}+h\) is prime with
\(p\nmid2D\), then some \(k\ge1\) satisfies
\[
 p\mid U_{2^{r+jk}}+h\qquad(j\ge1).
\]
These later numbers are larger than \(p\).
For every positive odd \(P\) and every integer \(h\),
the set of \(r\ge1\) with \(U_{2^r}+h\) composite has
positive lower natural density. Taking \(P=1\) proves
Saito's full Problem~1.8.
\end{theorem}
\begin{proof}
Let \(V_0=2,V_1=P,V_{n+2}=PV_{n+1}+V_n\), and put
\(f_r=U_{2^r}\), \(\ell_r=V_{2^r}\), \(r\ge1\).
For the roots \(\alpha,\beta\) of \(X^2-PX-1\),
\[
 U_n=\frac{\alpha^n-\beta^n}{\alpha-\beta},
 \qquad V_n=\alpha^n+\beta^n,\qquad \alpha\beta=-1.
\]
Expansion at even indices gives
\begin{equation}\label{exm:lucas-state}
 f_{r+1}=f_r\ell_r,\qquad
 \ell_{r+1}=\ell_r^2-2,\qquad
 \ell_r^2-Df_r^2=4.
\end{equation}
Here \(f_1=P\), \(\ell_1=P^2+2\), all \(f_r\) are odd,
and all \(\ell_r\equiv3\pmod4\).
Since \(P^2\equiv1\pmod8\),
\[
 v_2(\ell_1+1)=2,\qquad
 \ell_{r+1}+1=(\ell_r-1)(\ell_r+1),
\]
and consequently \(v_2(\ell_r+1)=r+1\). Also
\[
 f_{r+2}-f_r
 =f_r(\ell_r+1)(\ell_r^2-\ell_r-1).
\]
The first and last factors are odd; hence
\begin{equation}\label{exm:lucas-rate}
 v_2(f_{r+2}-f_r)=r+1.
\end{equation}

The two parity subsequences of \(f_r\) converge in
\(\mathbb Z_2\) to odd limits \(u_0,u_1\). In the
series for \(u_{r\bmod2}-f_r\), the first difference
has valuation \(r+1\), and all later ones have valuation
at least \(r+3\). Thus
\begin{equation}\label{exm:exact-lucas-limit}
 v_2(f_r-u_{r\bmod2})=r+1.
\end{equation}
Since \(\ell_r\to-1\), passing to the limit in
\eqref{exm:lucas-state} gives
\[
 u_1=-u_0,\qquad Du_i^2=-3.
\]
Neither limit is an ordinary integer, because \(D>0\).
For either limit \(u\) and every \(c\in\mathbb Z\),
\[
 v_2(u-c)+v_2(u+c)=v_2(Dc^2+3),
\]
by the product identity
\((u-c)(u+c)=-(Dc^2+3)/D\).
Both valuations are finite and nonnegative, so each
is at most the right-hand side.

Now suppose \(p=f_r+h\) is prime as in the theorem.
Apply this bound at \(c=1-h\) and \(c=-1-h\).
By \eqref{exm:exact-lucas-limit} and \(r\ge B_{P,h}+1\),
\begin{equation}\label{exm:two-budget}
 v_2(p-1)\le B_{P,h},\qquad v_2(p+1)\le B_{P,h}.
\end{equation}
Indeed, in \(f_r-c=(f_r-u)+(u-c)\), the first term
has strictly larger valuation.

Choose a root \(\gamma\) of \(X^2-PX-1\) in
\(\mathbb F_{p^2}\), with multiplicative order \(t\).
The roots are nonzero and distinct since \(p\nmid2D\).
If they lie in \(\mathbb F_p\), then \(t\mid p-1\).
Otherwise Frobenius interchanges the roots, whose product
is \(-1\), so \(\gamma^p=-\gamma^{-1}\) and
\(t\mid2(p+1)\). In either case \eqref{exm:two-budget}
implies \(v_2(t)\le B_{P,h}+1\le r\).
Thus \(z=\gamma^{2^r}\) has odd order \(q\).
Choose \(k\ge1\) with \(2^k\equiv1\pmod q\), taking
\(k=1\) when \(q=1\). Then
\[
 \gamma^{2^{r+jk}}=\gamma^{2^r}\qquad(j\ge0).
\]
The other root is \(\delta=-\gamma^{-1}\); because the
exponents are even, the same identity holds for \(\delta\).
Binet's formula is valid in \(\mathbb F_{p^2}\), so
\[
 U_{2^{r+jk}}\equiv U_{2^r}\pmod p.
\]
This proves the divisibility. Strict real increase of
\(U_{2^r}\) makes the divisor proper.

Choose \(R\) so large that for \(r\ge R\), the values
\(f_r+h\) exceed \(1\), exceed every prime divisor of
\(2D\), and satisfy \(r\ge B_{P,h}+1\).
If no such value is prime, all are composite and the
density is \(1\). Otherwise choose one prime value;
the return just proved gives a composite progression
of indices with density \(1/k>0\). No existence
assumption for a prime seed is required.
\end{proof}

The late-index condition matters. For \(P=1,h=4,r=2\),
one has \(F_4+4=7\). Modulo \(7\), the dyadic state is
\((f_2,\ell_2)=(3,0)\). One iteration makes
\(f_r\equiv0\pmod7\) forever, so this early prime never
divides a later shifted value. Here \(B_{1,4}=7\);
the theorem does not apply. In contrast \(h=-16\)
gives \(B_{1,-16}=3\) and \(F_{16}-16=971\), a valid
seed at \(r=4\). A root of \(X^2-X-1\) modulo \(971\)
is \(174\), of order \(485\), and
\(\operatorname{ord}_{485}(2)=48\). Therefore
\[
 971\mid F_{2^{\,4+48j}}-16\qquad(j\ge1).
\]
These exact examples illustrate the mechanism; the full
all-shift theorem rests on the proof, not on finite testing.

\subsection{What the five targets establish}\label{exm:status}

Problems~1.7 and~1.8 of \cite{SaitoIntervals} are answered
with their full original quantifiers by the two corpus proofs
reconstructed above. Their status here is complete proof
previously present in the supplied corpus, not a new
discovery in this paper. The square-clock irrationality
question, the Mills cubic-Pisot alternative, and the universal
pure cubic-clock composite question are not closed.
Their partial theorems identify an irrational derived series,
a necessary transcendental gap exponent, and a three-adic
exceptional set. None replaces the original question.


% Standalone mathematical fragment; no preamble or bibliography.
% Requires amsmath, amssymb, amsthm and theorem/lemma/corollary/example environments.
% Bibliography keys: FR2026, FR2016, Beukers2006, Lepetit2021.

\section{The external units conjecture for values of E-functions}
\label{ev:units-section}

\paragraph{The external problem.}
Fischler and Rivoal formulate the following problem as Conjecture~1.2
of \cite{FR2026}: if $\mathcal E$ is the ring of values of $E$-functions
at algebraic points, is
\[
 \mathcal E^\times
 =\{\alpha e^\beta:\alpha\in\overline{\mathbb Q}^{\times},
                         \ \beta\in\overline{\mathbb Q}\}?
\]
Their Theorem~1.3 proves the assertion for elements which are finite
algebraic-coefficient exponential sums.  We prove the conjecture on a
larger explicit class: all $E$-values lying in the field generated by
those exponential values and one exponential-integral value.  The
unrestricted conjecture is not concluded from this result.

Unless otherwise specified, $E$-functions have the strict arithmetic
growth bounds $C^{n+1}$ on coefficient conjugates and common
denominators, as in the main formulation of \cite{FR2026}.  The
arguments also apply to the larger ring obtained from Siegel's
original bounds $(n!)^\varepsilon$ for each $\varepsilon>0$; the
additional analytic input for that extension is specified below.

Put $A=\overline{\mathbb Q}$ and define
\[
 \mathcal B=A[e^\beta:\beta\in A],\qquad
 \mathcal L=\operatorname{Frac}(\mathcal B),\qquad
 \mathcal K=\operatorname{Frac}(\mathcal E).
\]
The ring $\mathcal B$ already contains every $e^{-\beta}$.
For $a\in A\setminus\{0\}$ write
\[
 \eta_a=\operatorname{Ein}(a),\qquad
 \operatorname{Ein}(z)=\sum_{n\ge1}\frac{(-1)^{n-1}z^n}{n\,n!}.
\]
No arithmetic assumption on Euler's constant is made in this section.

\subsection{A polynomial parameter for repeated winding}

We recall the analytic inputs and state exactly where they enter.
Choose a direction in $(-\pi/2,\pi/2)$ avoiding all anti-Stokes
directions of $E$-functions; this is possible because the family of
$E$-functions is countable and each has finitely many such directions.
For this fixed direction, the canonical decomposition
of an $E$-function has the form
\[
 f(z)=\sum_\rho f_\rho(z)e^{\rho z},\qquad \rho\in A,
\]
with finite support.  The components $e^{\rho z}f_\rho(z)$ solve an
$E$-operator and continue along all paths in $\mathbb C^\times$
\cite[Theorem~2.3 and the paragraph preceding it]{FR2026}.
This is a homogeneous annihilating operator for the full components,
even if a smaller inhomogeneous equation has also been used to
describe the original function.  Such an operator is regular singular
at zero with rational exponents
\cite[\S3.2, proof of Proposition~3, p.~41]{FR2016}.  Consequently its
local monodromy is quasi-unipotent.  Multiplication by the
single-valued, nowhere-zero factor $e^{-\rho z}$ identifies the
monodromy on $e^{\rho z}f_\rho$ with that on $f_\rho$.

For the broader Siegel definition, the same existence and local
structure follow from \cite[Proposition~16 and Theorem~7(a,b,d)]{Lepetit2021};
the continuation, canonical decomposition, and zero-division inputs
extend as explained in \cite[pp.~3--4]{FR2026}.  In particular, the
quasi-unipotence used below does not require the unproved equivalence
of the strict and broad definitions.

Let $M^j$ denote continuation around a based loop at $1$ of winding
number $j$ about zero.  Define
\[
 \Phi_j(f(1))(X)=\sum_\rho (M^j f_\rho)(1)e^{\rho X}.
\]
These are injective $A$-algebra homomorphisms from $\mathcal E$ to
\[
 \mathcal R=\left\{\sum_{\rho\in A}c_\rho e^{\rho X}:
                 c_\rho\in\mathbb C,\ \text{finite support}\right\},
\]
and satisfy $\Phi_j(\xi)(1)=\xi$ and
$\Phi_j(e^\beta)=e^{\beta X}$.
Here is the well-definedness argument, including its continuation
step.  If $f(1)=g(1)$, Beukers' theorem \cite{Beukers2006} gives an
$E$-function $h$ with $f-g=(z-1)h$.  Canonical uniqueness gives
$f_\rho-g_\rho=(z-1)h_\rho$.  After continuation, $h_\rho$ remains
holomorphic at $1$, so the corresponding continued component values
agree.  Componentwise addition and convolution for products give the
homomorphism property.  Continuing the sum recovers the entire,
single-valued function $f$, so evaluation at $X=1$ is a left inverse.

\begin{lemma}[Polynomial winding on a finite algebra]
\label{ev:polynomial-winding}
For every finitely generated $A$-subalgebra $S\subset\mathcal E$, there
are an integer $m\ge1$ and an injective $A$-algebra homomorphism
\[
 \Psi:S\longrightarrow\mathcal R[T]
\]
such that $\Psi(\xi)|_{T=j}=\Phi_{mj}(\xi)$ for every integer $j$.
It extends to the fraction field of $S$.
\end{lemma}
\begin{proof}
Choose representing $E$-functions for finitely many generators of $S$
and their finitely many canonical components.  A common positive
integer $m$ makes every relevant monodromy operator $M^m$ unipotent.
If $M^m=I+N$ with $N$ nilpotent, then
\[
 (M^m)^j=\sum_{k\ge0}\binom jk N^k
\]
is a finite sum and is polynomial in $j$, for all integers $j$.
Every coefficient of $\Phi_{mj}(\xi)$ for a chosen generator is
therefore polynomial in $j$.  Replace $j$ by an indeterminate $T$.

If a polynomial relation with coefficients in $A$ holds among the
generators, its images vanish after setting $T$ equal to every integer.
A polynomial over the domain $\mathcal R$ with infinitely many roots
is zero.  Hence the assignment is well-defined and multiplicative on
$S$.  Evaluation at $X=1$ gives the constant polynomial $\xi$ because
this holds at every integer $T$.  Thus the map is injective and extends
to fractions.  No common monodromy exponent for all of $\mathcal E$
is asserted or needed.
\end{proof}

For fixed nonzero algebraic $a$, the canonical components of
$\operatorname{Ein}(az)$, with a compatible branch, have the form
\[
 f_0(z)=C+\log z,
 \qquad
 f_{-a}(z)=e^{az}\bigl(\operatorname{Ein}(az)-C-\log z\bigr).
\]
The constant $C$ depends on the chosen branch.  The two phases are
$0$ and $-a$.  Continuing adds $2\pi i j$ to $f_0$ and subtracts
$2\pi i j e^{az}$ from $f_{-a}$.  Writing
\[
 c=e^a,\qquad B=c(\eta_a-C),\qquad \lambda=2\pi i m,
\]
the polynomial map of Lemma~\ref{ev:polynomial-winding}, on an algebra
containing $c$ and $\eta_a$, consequently satisfies
\begin{equation}\label{ev:eta-polynomial-image}
 \Psi(c)=e^{aX},\qquad
 \Psi(\eta_a)
 =C+Be^{-aX}+\lambda T(1-ce^{-aX}).
\end{equation}
The symbols $C,B,c$ on the right are complex coefficients.  Their use
does not assert that they are additional elements of the domain.

\begin{lemma}[Finite exponential phases with a parameter]
\label{ev:finite-phases}
Let $a_1,\ldots,a_r\in A$ be linearly independent over $\mathbb Q$,
and put $U_i=e^{a_iX}$.  If $W\in\operatorname{Frac}(\mathcal R[T])$
is algebraic over $\mathbb C(T,U_1,\ldots,U_r)$, then for some $q\ge1$,
\[
 W\in\mathbb C(T,e^{a_1X/q},\ldots,e^{a_rX/q}).
\]
Moreover,
\begin{equation}\label{ev:Laurent-intersection}
 \mathcal R[T]\cap\mathbb C(T,U_1,\ldots,U_r)
   =\mathbb C[T,U_1^{\pm1},\ldots,U_r^{\pm1}].
\end{equation}
\end{lemma}
\begin{proof}
Extend $a_1,\ldots,a_r$ to a rational basis for the span of the finitely
many phases occurring in a numerator and denominator for $W$, and
clear denominators.  The resulting exponential variables are
algebraically independent over $\mathbb C(T)$, by linear independence
of exponential functions with distinct phases.  The additional basis
variables give a purely transcendental extension and cannot contribute
to an element algebraic over the indicated smaller field.  This proves
the first assertion.

In the ring case the expression initially is a Laurent polynomial in
these variables with coefficients in $\mathbb C[T]$.  Independence
removes all monomials involving an additional basis variable.  Invariance
under the finite Kummer transformations of $e^{a_iX/q}$ then forces
each surviving exponent to be divisible by $q$.  This proves
\eqref{ev:Laurent-intersection}.
\end{proof}

We also use the elementary constant-extension fact that an irreducible
polynomial over $A(U_1,\ldots,U_r,V)$ remains irreducible over
$\mathbb C(U_1,\ldots,U_r,V)$.  After clearing denominators this is
absolute irreducibility over the algebraically closed field $A$.
It follows, for example, by expressing a hypothetical factorization
as a finite polynomial system and applying the weak Nullstellensatz.
In particular, adjoining complex constants preserves the degree of a
finite extension of such a rational function field.

\subsection{Complete unit classification in an explicit enlarged field}

\begin{theorem}[Exponential values with one Ein value]
\label{ev:one-Ein-field}
Let $a\in A\setminus\{0\}$ and put $F_a=\mathcal L(\eta_a)$.
Then:
\begin{enumerate}
\item $\eta_a$ is transcendental over $\mathcal L$.
\item $F_a$ is relatively algebraically closed in $\mathcal K$.
\item $\mathcal E\cap F_a=\mathcal B[\eta_a]$.
\item
\[
 \mathcal E^\times\cap F_a
 =\{\alpha e^\beta:\alpha\in A^\times,\ \beta\in A\}.
\]
\end{enumerate}
Thus the external units conjecture holds for every $E$-value that is
algebraic over $F_a$, as well as for every $E$-value in $F_a$ itself.
\end{theorem}
\begin{proof}
Every expression under consideration uses only finitely many exponential
values.  Include $a$ among their phases, choose a rational basis
beginning with $a$, and divide all its members by a common positive
integer $d$ that clears the coordinates of the finitely many phases.
The resulting basis $b_1,\ldots,b_r$ satisfies $a=db_1$ and every
phase under consideration is an integer combination of the $b_i$.
Adjoin the corresponding values $e^{b_i}$ and their inverses.
Apply Lemma~\ref{ev:polynomial-winding} to an algebra
containing the required $E$-values.  Put $U_i=e^{b_iX}$ and write
$V=\Psi(\eta_a)$.  Equation~\eqref{ev:eta-polynomial-image} becomes
\begin{equation}\label{ev:birational-map}
 V=C+B/U_1^d+\lambda T(1-c/U_1^d),\qquad c=e^a.
\end{equation}
Its coefficient of $T$ is nonzero.  Hence $U_1,\ldots,U_r,V$ are
algebraically independent over $\mathbb C$ and
\[
 \mathbb C(U_1,\ldots,U_r,V)
   =\mathbb C(U_1,\ldots,U_r,T).
\]
A polynomial relation over $A$ among the actual values
$e^{b_1},\ldots,e^{b_r},\eta_a$ would give one among these independent
images.  This proves the first assertion, including transcendence over
the union of all the finite exponential fields.

For the second assertion, let $v\in\mathcal K$ be algebraic over
$F_a$.  Its algebraic equation uses finitely many exponential values,
so enlarge the finite basis as above.  Include a numerator and a
denominator representing $v$ among the generators used to construct
$\Psi$.  Its image $W$ is algebraic over
$\mathbb C(T,U_1,\ldots,U_r)$.  Lemma~\ref{ev:finite-phases} gives
\[
 W\in\mathbb C(T,e^{b_1X/q},\ldots,e^{b_rX/q})
\]
for some $q\ge1$.  Adjoin the actual exponential roots $e^{b_i/q}$;
their monodromy is trivial, so the same construction extends to them.
Set
\[
 F_q=A(e^{b_1/q},\ldots,e^{b_r/q},\eta_a).
\]
The images of its $r+1$ generators are algebraically independent over
$\mathbb C$; solving \eqref{ev:birational-map} for $T$ shows that
$W\in\mathbb C\Psi(F_q)$.  The degree of $v$ over $F_q$ is preserved
under extension of constants from $A$ to $\mathbb C$.  That degree
is therefore one, and $v\in F_q\subset F_a$.

For the ring intersection, let $x\in\mathcal E\cap F_a$, and choose
a finite basis as before so that
\[
 x=R(e^{b_1},\ldots,e^{b_r},\eta_a),\qquad
 R\in A(U_1,\ldots,U_r,V).
\]
Write $R=P/Q$ with coprime $P,Q$ in the Laurent polynomial ring
\[
 A[U_1^{\pm1},\ldots,U_r^{\pm1},V].
\]
They remain coprime after extending constants to $\mathbb C$.
Let $\sigma$ be the substitution in \eqref{ev:birational-map}.
By \eqref{ev:Laurent-intersection},
\[
 \sigma(P)/\sigma(Q)=\Psi(x)
   \in\mathbb C[T,U_1^{\pm1},\ldots,U_r^{\pm1}].
\]
After inverting $U_1^d-c$, the substitution $\sigma$ is an isomorphism
of the two Laurent polynomial rings, because
\[
 T=\frac{U_1^d(V-C)-B}{\lambda(U_1^d-c)}.
\]
Consequently every common irreducible factor of $\sigma(P)$ and
$\sigma(Q)$ would have to divide $U_1^d-c$.

For each root $\zeta$ of $U_1^d-c$, at $U_1=\zeta$ the substituted value of $V$ is
$C+B/c=\eta_a$, independently of $T$.  A nonzero Laurent polynomial
over $A$ cannot vanish identically under
$(U_1,V)=(\zeta,\eta_a)$: the numbers $\zeta$ and $\eta_a$ are
algebraically independent over $A$, since $\zeta^d=c$ and the first
part of the proof makes $c,\eta_a$ algebraically independent.
Thus no factor $U_1-\zeta$ divides $\sigma(Q)$.
The substituted numerator and denominator are coprime.  Since their
quotient lies in the target ring, $\sigma(Q)$ must be a unit there,
of the form $d_0U_1^{n_1}\cdots U_r^{n_r}$ with $d_0\in\mathbb C^\times$.
Injectivity of $\sigma$ then gives the same identity for $Q$, and its
algebraic coefficients imply $d_0\in A^\times$.  The original
denominator is a unit.  Therefore
$x\in A[e^{\pm b_1},\ldots,e^{\pm b_r},\eta_a]\subset\mathcal B[\eta_a]$.
The reverse inclusion is immediate.

Finally, if $x\in\mathcal E^\times\cap F_a$, both $x$ and $x^{-1}$
lie in $\mathcal B[\eta_a]$.  Because $\eta_a$ is transcendental over
$\mathcal L$, the units of this polynomial ring are exactly the units
of $\mathcal B$.  A unit and its inverse in $\mathcal B$ involve only
finitely many phases.  Clearing their rational coordinates identifies
them with mutually inverse Laurent polynomials in algebraically
independent exponential values.  Such units are precisely monomials,
so $x=\alpha e^\beta$ with $\alpha\in A^\times$ and $\beta\in A$.
Their inverses plainly belong to $\mathcal E$.  The final sentence of
the theorem follows from relative algebraic closure.
\end{proof}

\begin{corollary}[Explicit non-units and non-members]
\label{ev:explicit-nonunits}
Let $P\in\mathcal B[Y]$ be nonconstant.  Then $P(\eta_a)$ is a
nonzero non-unit of $\mathcal E$, and
\[
 \frac1{P(\eta_a)}\notin\mathcal E.
\]
More generally, a rational function in $\eta_a$ with coefficients in
$\mathcal L$ is an $E$-value exactly when it belongs to
$\mathcal B[\eta_a]$.
\end{corollary}
\begin{proof}
The first assertion follows from the transcendence of $\eta_a$ over
$\mathcal L$ and the unit classification.  The second is the proved
ring intersection.
\end{proof}

\begin{example}
Writing $\eta=\operatorname{Ein}(1)$, the numbers
\[
 \eta^2+e\eta+e^{\sqrt2},\qquad
 e^{\sqrt3}\eta+e^{-\sqrt2},\qquad
 \eta-e^a\quad(a\in A)
\]
are nonzero $E$-values whose reciprocals are not $E$-values.  These
examples use the same complete classification; they are not additional
external conjectures counted as separately solved.
\end{example}

\subsection{Several independent Ein scales}

The complementary family below allows several exponential-integral
values without adjoining all exponential values to their base field.
This distinction is essential: the two field constructions are not
silently combined.

\begin{theorem}[Several Ein values]
\label{ev:several-Ein}
Let $a_1,\ldots,a_r\in A\setminus\{0\}$ be linearly independent over
$\mathbb Q$, with no reality or positivity restriction.  Put
\[
 F=A(\eta_{a_1},\ldots,\eta_{a_r}),\qquad
 D=A[\eta_{a_1},\ldots,\eta_{a_r}].
\]
The displayed generators are algebraically independent, $F$ is
relatively algebraically closed in $\mathcal K$, and
\[
 \mathcal E\cap F=D,
 \qquad
 \mathcal E^\times\cap F=A^\times.
\]
\end{theorem}
\begin{proof}
For a fixed common lateral direction, write
\[
 \Phi_j(\eta_{a_i})=C_i+j\tau+(B_i-j\tau e^{a_i})e^{-a_iX},
 \qquad \tau=2\pi i.
\]
There is at most one integer $j$ making any specified exponential
coefficient zero.  Choose two distinct integers $j_0,j_1$ avoiding all
these finitely many exceptions.  In either image the $r$ generators
are affine functions of $r$ algebraically independent exponential
variables.  This proves their algebraic independence over $A$.

Let $v\in\mathcal K$ be algebraic over $F$, and identify the generators
abstractly with $U_1,\ldots,U_r$.  In either image, finite-phase descent
places the resulting extension, after extending constants to
$\mathbb C$, inside a translated Kummer extension
\[
 \mathbb C(Z_1,\ldots,Z_r)/\mathbb C(U_1,\ldots,U_r),
 \qquad
 U_i=C_i+j\tau+(B_i-j\tau e^{a_i})Z_i^q.
\]
Its group is $(\mathbb Z/q\mathbb Z)^r$.  The inertia at the generic
point of $U_i=C_i+j\tau$ is its $i$th coordinate subgroup, and these
inertia groups generate the entire group.  Every nontrivial
intermediate extension is therefore ramified along at least one of
these finite hyperplanes.

The two realizations describe the same abstract extension of
$\mathbb C(U_1,\ldots,U_r)$, because extension of constants preserves
its degree and defining polynomial.  The second Kummer realization
is unramified at the generic point of every finite hyperplane from
the first: for the same coordinate the constants differ by
$(j_1-j_0)\tau$, and hyperplanes in different coordinates are distinct.
Thus all coordinate inertia groups of the first realization act
trivially on the intermediate field.  That field is the base field,
so the original degree is one and $v\in F$.

For the ring assertion, write an element of $F\cap\mathcal E$ as
$P(\eta_{a_1},\ldots,\eta_{a_r})/
 Q(\eta_{a_1},\ldots,\eta_{a_r})$ with coprime $P,Q\in A[U_1,\ldots,U_r]$.
Finite-phase descent in the ring case implies that its image is a
Laurent polynomial in the variables $e^{-a_iX}$.  Its denominator
can therefore have irreducible factors only among
$U_i-(C_i+j\tau)$.  Applying this conclusion for both selected
windings leaves no possible factor: the two lists of hyperplanes
are disjoint.  Hence $Q$ is constant and the intersection is $D$.
If an element of $F$ is a unit of $\mathcal E$, it and its reciprocal
belong to the polynomial ring $D$.  Its only units are $A^\times$.
\end{proof}

\paragraph{What remains in the external problem.}
Theorems~\ref{ev:one-Ein-field} and \ref{ev:several-Ein} completely
classify units in their specified classes.  An arbitrary $E$-value
need not be algebraic over any of the fields used here.  No reduction
of every unit of $\mathcal E$ to those fields has been proved.  That
missing global step is precisely why Conjecture~1.2 is recorded as a
partial solution in the ten-problem register.


% Integration fragment. Main external slot: the nondecreasing variant of #270.
% Auxiliary application: a regular-growth obstruction related to #265.
% Bibliography keys are listed in series_external_notes.md.

\section{The nondecreasing variant of Erd\H{o}s problem 270}
\label{es:270}

The external question is whether
\begin{equation}\label{es:270-series}
 S(f)=\sum_{n\ge1}\frac{n!}{(n+f(n))!}
\end{equation}
is irrational whenever $f:\mathbb N\to\mathbb N$ is nondecreasing
and tends to infinity.  Crmari\'c and Kova\v{c}
\cite{CrKo2025} disprove the version without monotonicity: its
values fill $(0,\infty)$.  They explicitly retain the monotone
question and prove that its set of possible values has measure zero.
We do not close the monotone question.  We prove a concrete
subfamily, with a bounded-numerator strengthening.

\paragraph{A necessary historical correction.}
The 2025 paper also singled out $f(n)=n$.  That example must not be
counted as a new solution here.  Lambrinoudis's primary repository
\cite{Lambrinoudis270} credits Crmari\'c--Kova\v{c}'s elementary
proof, posted by Kova\v{c} in July 2026, and explains that its
transcendence follows from an earlier Siegel--Shidlovskii argument.
The same repository contains broader affine-family claims.  The
results below use neither those claims nor their analytic inputs.

\begin{lemma}[Clearing denominators in a short window]
\label{es:window}
Let $f:\mathbb N\to\mathbb N$ be nondecreasing, and write
\[
 g_n=n+f(n),\qquad D_n=\frac{g_n!}{n!},\qquad
 r_n=\frac{D_{n+1}}{D_n}.
\]
Suppose that, for all sufficiently large $N$, an integer
$h_N\in\{1,\ldots,N-1\}$ satisfies
\begin{equation}\label{es:window-gap}
 g_N-g_{N-h_N}\ge N.
\end{equation}
Then the integer
\[
 L_N=\frac{g_N!}{(N-h_N)!}
     =D_N H_N,\qquad
 H_N=\frac{N!}{(N-h_N)!},
\]
is divisible by every $D_k$ with $1\le k\le N$.
If, in addition,
\begin{equation}\label{es:window-small}
 \frac{H_N}{r_N}\longrightarrow0,
\end{equation}
then, for every bounded integer sequence $(u_n)$ having infinitely
many nonzero terms, the convergent series
\[
 \sum_{n\ge1}\frac{u_n}{D_n}
\]
is irrational.
\end{lemma}

\begin{proof}
If $k\le N-h_N$, then $g_N-g_k\ge N$.  The numerator
$g_N!/g_k!$ contains a product of $N$ consecutive positive integers,
and such a product is divisible by $N!$.  Consequently
\[
 \frac{D_N}{D_k}=\frac{g_N!}{g_k!}\frac{k!}{N!}
 \in\mathbb Z.
\]
If $N-h_N<k\le N$, then instead
\[
 \frac{L_N}{D_k}
   =\frac{g_N!}{g_k!}\frac{k!}{(N-h_N)!}
   \in\mathbb Z.
\]
This proves the divisibility assertion.

Since $H_N\ge1$, condition \eqref{es:window-small} implies
$r_N\to\infty$.  Hence all the series in question converge
absolutely.  Put $B=\sup_n|u_n|$, which is a positive integer, and
choose an arbitrarily large $N$ with $u_{N+1}\ne0$.  Since the
successive ratios tend to infinity, for all sufficiently large such
$N$ the tail satisfies
\[
 \frac{1}{2D_{N+1}}
 \le\left|\sum_{k>N}\frac{u_k}{D_k}\right|
 \le\frac{2B}{D_{N+1}}.
\]
For the lower bound, the first nonzero term has magnitude at least
$1/D_{N+1}$, while the sum of subsequent magnitudes is
$o(1/D_{N+1})$.  The upper bound follows from a geometric majorant.

If the whole sum were $a/b\in\mathbb Q$, with $b>0$, then
\[
 bL_N\left(\frac ab-\sum_{k\le N}\frac{u_k}{D_k}\right)
\]
would be a nonzero integer whose absolute value is at most
$2bB H_N/r_N\to0$.  This is impossible.
\end{proof}

\begin{theorem}[A power-growth subfamily of the monotone question]
\label{es:power270}
Let $c>0$, and let $f:\mathbb N\to\mathbb N$ agree eventually with
\[
 f(n)=\lfloor c n^\alpha\rfloor.
\]
The series \eqref{es:270-series} is irrational in each of the
following ranges:
\begin{equation}\label{es:power-range}
 \alpha>\frac32;
 \qquad\text{or}\qquad
 \alpha=\frac32,\quad c>\left(\frac23\right)^{3/2}.
\end{equation}
More strongly, replacing the numerator $1/D_n$ by $u_n/D_n$ for
any bounded integer sequence with infinitely many nonzero terms
preserves irrationality.
\end{theorem}

\begin{proof}
Changing finitely many terms changes the sum by a rational number,
so we may begin where $f$ is positive and nondecreasing.  With the
notation of Lemma~\ref{es:window},
\[
 g_n\sim c n^\alpha,
 \qquad
 g_{n+1}-g_n=(c\alpha+o(1))n^{\alpha-1}.
\]
The exact quotient is
\[
 r_n=\frac{\prod_{j=g_n+1}^{g_{n+1}}j}{n+1}.
\]
All factors in this numerator are asymptotic to $c n^\alpha$;
therefore
\begin{equation}\label{es:ratio-asymptotic}
 \log r_n=(c\alpha^2+o(1))n^{\alpha-1}\log n.
\end{equation}

First suppose $3/2\le\alpha<2$.  Choose
$t>1/(c\alpha)$ and put $h_N=\lceil tN^{2-\alpha}\rceil$.
Taylor's formula and the bounded floor error give
\[
 g_N-g_{N-h_N}=(c\alpha t+o(1))N>N
\]
for large $N$, so \eqref{es:window-gap} holds.  Moreover
\begin{equation}\label{es:height-upper}
 \log H_N\le h_N\log N
       =(t+o(1))N^{2-\alpha}\log N.
\end{equation}
When $\alpha>3/2$, comparison with
\eqref{es:ratio-asymptotic} proves $H_N/r_N\to0$.
At $\alpha=3/2$, the hypothesis on $c$ is exactly the condition
$1/(c\alpha)<c\alpha^2$.  We can then choose $t$ strictly between
these two numbers, and the same comparison again proves the limit.

For $\alpha=2$, take a fixed integer $h>1/(2c)$.
Then $g_N-g_{N-h}\sim2chN>N$, whereas
$\log H_N=O(\log N)=o(\log r_N)$.
For $\alpha>2$, take $h=1$; the same conclusions follow directly.
Lemma~\ref{es:window} completes all cases.
\end{proof}

\begin{corollary}[Polynomial factorial shifts]
\label{es:polynomial270}
Let $f$ be a polynomial of degree at least two with positive leading
coefficient and integer values at integers.  Begin the series where
$f(n)>0$.  Then
\[
 \sum_n\frac{n!}{(n+f(n))!}
\]
is irrational, as is every version with bounded integer numerators
having infinitely many nonzero terms.
\end{corollary}

\begin{proof}
If $f(n)=cn^d+O(n^{d-1})$, then
$g_{n+1}-g_n=cd n^{d-1}+O(n^{d-2})$.
For $d=2$ choose a fixed integer $h>1/(2c)$; for $d>2$ choose
$h=1$.  Condition \eqref{es:window-gap} follows, while
$\log H_N=O(\log N)$ and
$\log r_N=(cd^2+o(1))N^{d-1}\log N$.
Apply Lemma~\ref{es:window}.
\end{proof}

\begin{corollary}[Disjoint subseries]
\label{es:subseries270}
Under either Theorem~\ref{es:power270} or
Corollary~\ref{es:polynomial270}, let $A_1,\ldots,A_q$ be pairwise
disjoint infinite subsets of $\mathbb N$.  Then the numbers
\[
 1,\quad \sum_{n\in A_1}\frac1{D_n},\quad\ldots,\quad
 \sum_{n\in A_q}\frac1{D_n}
\]
are linearly independent over $\mathbb Q$.
\end{corollary}
\begin{proof}
An integer linear combination of the subseries with at least one
nonzero coefficient is a series with bounded integer numerators
and infinitely many nonzero terms.  The preceding results make
that combination irrational, so it cannot equal a rational constant.
\end{proof}

\paragraph{Relation to the original tail method.}
For $f(n)=n$, one has $D_n=(2n)!/n!$ and
$D_{n+1}/D_n=4n+2\in\mathbb Z$.  If its sum were rational, the
positive integral multiple of its $N$th tail would be bounded by
$b/(4N+1)\to0$.  This recovers the already known elementary proof.
For general monotone $f$, those successive quotients need not be
integers.  Lemma~\ref{es:window} supplies a different common
denominator and pays the explicit factor $H_N$.  The remaining
barrier is to make that factor smaller than $r_N$ for all monotone
$f\to\infty$; the present proof does not do so.  The constants in
\eqref{es:power-range} describe this argument's range and are not
claimed to be optimal boundaries of irrationality.



\section{The Meijer identity for \texorpdfstring{$\gamma+\log z$}{gamma + log z}}
\label{cmf:extension}

The ninth target is the analytic identity requested in
\cite[Section 9.2(1)]{CMF2026}. That paragraph also asks about
recurrence comparisons; the result here resolves its identity
component, with all ranks and all arguments on the principal branch.
It does not settle the broader dynamical conjecture treated in the
next section.

For $p\ge2$, put
\[
 G_p(z)=G^{p+1,p}_{p,p+1}
 \!\left(\begin{matrix}0,\ldots,0\\0,\ldots,0\end{matrix}
 \middle|z\right),\qquad \theta=z\frac{d}{dz}.
\]
We work first with $z>0$ and then continue to
$\mathscr U=\mathbb C\setminus(-\infty,0]$, using $\Log z$.
The Barnes convention is the clockwise loop around nonnegative
integer poles \cite[Section 16.17]{NISTDLMF}:
\begin{equation}\label{cmf:Barnes}
 G_p(z)=\frac1{2\pi i}\int_L\Phi_p(s)z^s\,ds,
 \qquad \Phi_p(s)=\Gamma(-s)^{p+1}\Gamma(1+s)^p.
\end{equation}

\begin{theorem}[The full argument-dependent identity]
\label{cmf:log-identity}
Define the differential operators
\begin{align*}
 \mathcal N_p(z)&=z(1+\theta)^{p-2}+(-1)^p\theta^{p-1},\\
 \mathcal D_p(z)&=z(1+\theta)^{p-1}+(-1)^p\theta^p.
\end{align*}
In these expressions $z$ means multiplication on the left.
For every $z\in\mathscr U$,
\begin{equation}\label{cmf:main-identity}
 \mathcal N_p(z)G_p(z)=\gamma+\Log z,
 \qquad \mathcal D_p(z)G_p(z)=1.
\end{equation}
At $z=1$ these recover the operators of
\cite[Proposition 1]{CMF2026}.
\end{theorem}

\begin{proof}[First proof: differential equation and its boundary constants]
The standard Meijer equation \cite[Section 16.21]{NISTDLMF} is
\[
 \bigl[z(1+\theta)^p+(-1)^p\theta^{p+1}\bigr]G_p=0.
\]
The operator identity $\theta z=z(\theta+1)$ gives
\[
 \theta\mathcal N_p=\mathcal D_p,
 \qquad
 \theta\mathcal D_p=z(1+\theta)^p+(-1)^p\theta^{p+1}.
\]
Therefore $\mathcal D_pG_p$ is constant, and
$\mathcal N_pG_p$ is affine in $\Log z$.

To determine the constants, the expansion at the Barnes pole zero is
\[
 \Phi_p(s)=(-1)^{p+1}s^{-p-1}(1+\gamma s+O(s^2)).
\]
Clockwise orientation contributes the negative residue. As
$z\to0^+$ this yields
\[
 G_p(z)=(-1)^p\left\{
   \frac{(\log z)^p}{p!}
   +\gamma\frac{(\log z)^{p-1}}{(p-1)!}
   +P_{p-2}(\log z)\right\}
   +O\bigl(z(1+|\log z|^p)\bigr),
\]
where $P_{p-2}$ is a polynomial of degree at most $p-2$.
The same remainder estimate holds after any fixed number of
Euler derivatives. One may see this directly from the convergent
residue expansion: its $n$th summand is $z^n$ times a polynomial
in $\log z$ of degree at most $p$, with coefficients bounded by
$O((1+\log(n+1))^p/n!)$ up to fixed polynomial factors in $n$.

It follows that $\mathcal D_pG_p\to1$ and
$\mathcal N_pG_p-\log z\to\gamma$. These determine both
constants. Analytic continuation gives the result on $\mathscr U$.
\end{proof}

\subsection{A residue hierarchy beyond the first two identities}

The same mechanism produces higher polynomials in $\gamma+\Log z$
and zeta values. This extension also supplies an independent
telescoping verification of Theorem~\ref{cmf:log-identity}.

\begin{lemma}[Weighted Barnes telescoping]
\label{cmf:telescoping}
For every polynomial $Q\in\mathbb C[s]$, define
\[
 \mathcal T_Q=zQ(\theta)+(-1)^p\theta Q(\theta-1).
\]
Then
\begin{equation}\label{cmf:residue-identity}
 \mathcal T_QG_p(z)
 =z\operatorname{Res}_{s=-1}\bigl(Q(s)\Phi_p(s)z^s\bigr).
\end{equation}
\end{lemma}
\begin{proof}
The functional equation of $\Gamma$ gives
$\Phi_p(s-1)=(-1)^{p+1}s\Phi_p(s)$. If
$F_z(s)=Q(s)\Phi_p(s)z^s$, the Barnes integrand for
$\mathcal T_QG_p$ is consequently $z(F_z(s)-F_z(s-1))$.
Summing its negative residues at $s=0,\ldots,N$ telescopes to
\[
 z\operatorname{Res}_{s=-1}F_z(s)
 -z\operatorname{Res}_{s=N}F_z(s).
\]
The last residue tends to zero locally uniformly on $\mathscr U$:
reflection gives
$\Phi_p(s)=(-\pi)^p\Gamma(-s)\csc^p(\pi s)$, so its
coefficients at $s=N$ are bounded by a fixed power of
$1+\log(N+1)$ times $1/N!$. Multiplication by $Q(s)z^s$
adds only polynomial factors in $N$ and the factor $|z|^N$ on
each compact set. This also justifies termwise application of
the finite differential operator, proving the identity.
\end{proof}

\begin{theorem}[A zeta-polynomial hierarchy]
\label{cmf:hierarchy}
For $0\le j\le p-1$, set
\[
 \mathcal T_{p,j}=z(1+\theta)^{p-1-j}+(-1)^p\theta^{p-j},
 \qquad T=\gamma+\Log z.
\]
Then
\begin{equation}\label{cmf:hierarchy-main}
 \mathcal T_{p,j}G_p(z)
 =[t^j]\exp\!\left(
 Tt+\sum_{m\ge2}\frac{1+p(1+(-1)^m)}{m}\zeta(m)t^m
 \right).
\end{equation}
The first four resulting expressions, when their indices are
available, are
\[
 1,\quad T,\quad
 \frac{T^2+(2p+1)\zeta(2)}2,\quad
 \frac{T^3+3(2p+1)\zeta(2)T+2\zeta(3)}6.
\]
\end{theorem}
\begin{proof}
Use Lemma~\ref{cmf:telescoping} with
$Q(s)=(s+1)^{p-1-j}$. At $s=-1+t$,
\[
 \Phi_p(-1+t)=t^{-p}\Gamma(1-t)
                 \left(\frac{\pi t}{\sin\pi t}\right)^p.
\]
Thus the right side of \eqref{cmf:residue-identity} is
\[
 [t^j]e^{t\Log z}\Gamma(1-t)
                 \left(\frac{\pi t}{\sin\pi t}\right)^p.
\]
The convergent local expansions
\[
 \log\Gamma(1-t)=\gamma t+\sum_{m\ge2}\frac{\zeta(m)}m t^m,
 \quad
 \log\frac{\pi t}{\sin\pi t}
   =\sum_{m\ge2}\frac{1+(-1)^m}{m}\zeta(m)t^m
\]
give \eqref{cmf:hierarchy-main}. Only finitely many coefficients
are used for each $j$.
\end{proof}

\begin{corollary}[Exact kernel of weighted telescoping]
\label{cmf:kernel}
For fixed $p$, the map $Q\mapsto\mathcal T_QG_p$ has image
$\mathbb C[\gamma+\Log z]_{\le p-1}$ and kernel
$((s+1)^p)\subset\mathbb C[s]$. More explicitly, if $H_j(T)$
denotes the coefficient polynomial in
\eqref{cmf:hierarchy-main}, then
\[
 \mathcal T_QG_p(z)=
 \sum_{\ell=0}^{p-1}\frac{Q^{(\ell)}(-1)}{\ell!}
 H_{p-1-\ell}(\gamma+\Log z).
\]
\end{corollary}
\begin{proof}
Taylor-expand $Q(-1+t)$ in the residue formula. The residue uses
only its terms of degrees at most $p-1$. The polynomials
$H_0,\ldots,H_{p-1}$ have distinct degrees and leading
coefficients $1/j!$, so they form a basis of the stated image.
Their linear independence identifies the kernel exactly.
\end{proof}

\begin{remark}[Scope of the resolution]
The $j=0,1$ cases are precisely the denominator and numerator
identities, now for arbitrary $z$. The factor $z$ on the left is
essential. Substituting arbitrary $z$ into the operators originally
written only for evaluation at $1$ would miss it. These are exact
analytic identities. They do not establish the irrationality of
$\gamma$, nor an integer-denominator estimate for any recurrence
derived from them. No claim of historical priority for identities
equivalent to known special-function formulas is made.
\end{remark}


\section{The all-trajectory CMF conjecture and an all-rank scalar theorem}
\label{cmf:convergence}

The tenth target is Conjecture~1 of \cite{CMF2026}. For rank
$r\ge3$, its trajectories are
\[
 \mathbf t=(a_1,\ldots,a_{r-1};b_1,\ldots,b_r),\quad
 a_i\le0,\quad b_j\ge1,
\]
with the smallest $b_j$ occurring at least twice. The conjecture
asks for generic dual-state convergence along every such trajectory.
The source proves the rank-three subfamily $b_1=b_2=b_3=1$.
We do not prove its all-rank, all-trajectory assertion. We establish
below a directly analyzable scalar approximation family and explain
the precise extra step needed to connect it to that assertion.

\subsection{A positive binomial family at every rank}

For an integer $r\ge2$, a fixed real $z>0$, and $0\le m\le n$, set
\[
 w_{n,m}(z)=\binom nm^r(n-m)!\,z^m,
 \qquad A_{n,m}=rH_m-(r-1)H_{n-m},\quad H_0=0.
\]
Define
\begin{equation}\label{cmf:scalar-approx}
 R_{r,n}(z)=
 \frac{\sum_{m=0}^n w_{n,m}(z)A_{n,m}}
      {\sum_{m=0}^n w_{n,m}(z)}.
\end{equation}
For positive rational $z$, every displayed approximant is rational.
At $r=3,z=1$, the harmonic/binomial structure agrees with the
scalar projection analyzed in \cite{Integrated2026}. The following
proof treats the whole family on its own terms; it requires no
unproved matrix-asymptotic hypothesis.

\begin{theorem}[Scalar convergence for every rank]
\label{cmf:scalar-convergence}
For fixed $r\ge2$ and $z>0$,
\begin{equation}\label{cmf:scalar-rate}
 R_{r,n}(z)=\gamma+\log z
       +O_{r,z}\bigl(n^{-(r-1)/(4r)}\bigr).
\end{equation}
The exponent is a convenient explicit bound from a concentration
argument; no optimality is asserted.
\end{theorem}
\begin{proof}
Let $x=x_n$ be the unique solution in $(0,n)$ of
\begin{equation}\label{cmf:saddle}
 x^r=z(n-x)^{r-1}.
\end{equation}
Uniqueness follows from strict increase of the left side divided by
$(n-x)^{r-1}$. This equation implies $x/n\to0$, and then
\[
 x\sim z^{1/r}n^{(r-1)/r}.
\]
Successive weights have ratio
\begin{equation}\label{cmf:weight-ratio}
 \rho(m)=\frac{w_{n,m+1}}{w_{n,m}}
        =\frac{z(n-m)^{r-1}}{(m+1)^r}.
\end{equation}
This is strictly decreasing as a real function on $(-1,n)$.
At $x-1$ its value exceeds $1$, and at $x$ it is less than $1$,
by \eqref{cmf:saddle}. Hence the index of a maximal weight differs
from $x$ by a bounded amount.

Put $s=x^{3/4}$. For all large $n$, the interval
$[x-2s,x+2s]$ lies in $[x/2,2x]\subset(0,n)$, and there
\[
 \frac{d}{dm}\log\rho(m)
 =-\frac{r-1}{n-m}-\frac r{m+1}\le-\frac{c_r}{x}
\]
for a fixed $c_r>0$. Summing log-ratios from a maximal index
to either boundary $|m-x|=s+O(1)$ gives
\[
 \frac{w_{n,m}}{\max_j w_{n,j}}
 \le \exp(-c'_r s^2/x)=\exp(-c'_r\sqrt x)
\]
at those boundaries, with a smaller positive constant if necessary.
Unimodality bounds every more distant weight by its nearer boundary
weight. Since there are at most $n+1$ terms, the probability measure
proportional to $w_{n,m}$ assigns the complement of
$|m-x|\le s$ mass at most
$O(n\exp(-c'_r\sqrt x))$.

Uniformly inside this window, $m=x(1+O(x^{-1/4}))$ and
$n-m=(n-x)(1+O(x^{3/4}/n))$. The elementary estimate
$H_k=\log k+\gamma+O(1/k)$ gives
\begin{align*}
 A_{n,m}
 &=\gamma+r\log m-(r-1)\log(n-m)+O(1/x)\\
 &=\gamma+r\log x-(r-1)\log(n-x)+O(x^{-1/4})\\
 &=\gamma+\log z+O(x^{-1/4}).
\end{align*}
On the complement, $|A_{n,m}|\le(2r-1)H_n=O_r(\log n)$.
Its total contribution is therefore
$O_{r,z}(n\log n\exp(-c'_r\sqrt x))$, negligible compared
with $x^{-1/4}$. Averaging proves
\eqref{cmf:scalar-rate}.
\end{proof}

\subsection{What this proves about the external conjecture}

Theorem~\ref{cmf:scalar-convergence} is an all-rank scalar result.
It does not identify the weighted expression with the dual matrix
output for every trajectory and every generic initial vector.
Such an identification, together with control of all competing
solution modes, would be required for Conjecture~1. The mere
existence of one convergent scalar family does not supply it.
Likewise, the convergence rate in \eqref{cmf:scalar-rate} is not
an irrationality measure: multiplying by a large common denominator
can make its integer linear form grow.

One elementary extension of the source's proved matrix family can
be recorded without this missing step.

\begin{proposition}[Integer dilation of a convergent trajectory]
\label{cmf:dilation}
Suppose a conservative matrix field has invertible transport
matrices $T(n\mathbf t)$ from the origin and
\[
 c_n T(n\mathbf t)^{-T}v\longrightarrow S
\]
for a specified generic vector $v$ and nonzero scalars $c_n$.
For every positive integer $d$, the same state-isolation conclusion
holds on the dilated trajectory $d\mathbf t$ with the same vector.
In particular, the source's proved rank-three family extends to
\[
 (da_1,da_2;d,d,d),\qquad a_1,a_2\le0,
\]
whenever the relevant transports are defined.
\end{proposition}
\begin{proof}
Conservativity makes the transport depend only on the endpoint.
Transport to $n(d\mathbf t)$ is therefore $T(dn\mathbf t)$.
The asserted convergence is the subsequence with indices $dn$ of
the assumed convergence. This proves the claim for every vector
for which the original limit holds.
\end{proof}

This is a consequence of the source theorem, not a new solution of
Conjecture~1. For example, a trajectory with unequal lower speeds
or higher rank is not covered by merely dilating the proved family.
The all-trajectory question remains unresolved in this paper.


\section{What has been closed, and the remaining proof obligations}
\label{ext:outcomes}

Saito's non-reversible recurrence and all-shift Fibonacci problems
receive full answers reconstructed from the corpus. The Meijer
extension resolves the selected analytic identity. The enlarged
$E$-value field is completely classified, but arbitrary units remain
outside that classification; likewise, the factorial-series theorem
does not cover every nondecreasing $f\to\infty$.

For the original gamma/Mills goals, the remaining obligations can be
stated without numerical ambiguity:
\begin{enumerate}
\item For $\gamma$, obtain nonzero \emph{primitive integer} linear
forms tending to zero, or independently establish the radical
membership in Corollary~\ref{ext:gamma-radical}. Neither follows
from the available analytic identities alone.
\item For $M$, exclude the actual cubic-Pisot alternative, including
its surviving three-adic trace limits $\lambda=\pm1$. Conditional
gap restrictions do not exclude a tower merely by describing it.
\item For $\xi_2$, connect rationality of the base to a contradiction
in its prime sequence. Irrationality of a reciprocal series that is
also irrational at rational bases cannot make that connection.
\item For the full CMF conjecture, control all competing modes of
the dual system on every admissible trajectory. One explicitly
convergent scalar expression does not identify a generic vector limit.
\end{enumerate}

The register therefore records two full answers, one resolved identity
component, and seven unresolved targets with partial results. RH,
$\gamma$ irrationality, and unconditional Mills transcendence are not proved.

\subsection{Verification record and reproducibility}

Separate internal proof checks covered the modular preperiods,
prime-return threshold, finite-phase descent, power-growth constant,
Barnes orientation, and scalar saddle. The phase argument explicitly
clears rational basis denominators, including values such as $e^{a/2}$.

The exact verifier passed 9,200 lattice cases, 23,091 denominator
divisibilities, and the displayed Lucas examples.
The numerical Meijer evaluations at 35 decimal working precision used
$(p,z)=(2,1/2),(2,2),(3,7/10)$; the largest reported residual among
the two main identities and the available $j=2$ identity was below
$10^{-30}$. Those decimal diagnostics do not replace the analytic proofs.
Appendix~\ref{app:verify} supplies the exact verifier. The TEX needs no
external input fragments, bibliography database, figures, or private data.



\clearpage
\appendix
\section{A compact exact-arithmetic verifier}
\label{app:verify}

The following standard-library Python program checks the displayed
finite lattice, common-denominator, and Lucas examples. Save the listing
as a UTF-8 file and run it with Python 3. It requires no network access.
Its finite output supplements the proofs and makes no claim to decide
an infinite conjecture by testing a bounded range.

\begin{lstlisting}[style=certificate]


"""Finite diagnostics for the ten-target paper; standard Python only.
The infinite theorems are proved in the manuscript, not by this script.
"""
from math import factorial, gcd, isqrt
from random import Random


def check_seeds():
    count = 0
    for D in range(1, 17):
        for C in range(D):
            if gcd(C, D) != 1:
                continue
            for p in range(-7, 8):
                for q in range(1, 13):
                    if gcd(p, q) != 1:
                        continue
                    h = gcd(D*p-C*q, D*q)
                    assert D*D % h == 0
                    a, b = (D*p-C*q)//h, D*q//h
                    k = D*D//h
                    assert gcd(a, b) == 1
                    assert (D*a+C*b, D*b) == (k*p, k*q)
                    assert gcd(D*a+C*b, D*b) == k
                    if gcd(q, D) == 1:
                        assert h == 1
                    count += 1
    return count


def check_windows():
    rng = Random(20261105)
    count = 0
    for trial in range(50):
        f = [1]
        for n in range(1, 31):
            f.append(f[-1] + rng.randrange(0, 5))
        g = [n + f[n] for n in range(31)]
        for N in range(3, 31):
            options = [h for h in range(1, N)
                       if g[N]-g[N-h] >= N]
            if not options:
                continue
            h = min(options)
            L = factorial(g[N]) // factorial(N-h)
            for k in range(1, N+1):
                Dk = factorial(g[k]) // factorial(k)
                assert L % Dk == 0
                count += 1
    return count


def v2(n):
    n = abs(n)
    assert n != 0
    return (n & -n).bit_length()-1


def check_lucas():
    for P in (1, 3, 5):
        f, ell = P, P*P+2
        for r in range(1, 10):
            assert ell*ell-(P*P+4)*f*f == 4
            assert v2(ell+1) == r+1
            assert v2(f*ell*(ell*ell-2)-f) == r+1
            f, ell = f*ell, ell*ell-2
    assert all(971 % d for d in range(2, isqrt(971)+1))
    assert (174*174-174-1) % 971 == 0
    assert pow(174, 485, 971) == 1
    assert all(pow(174, 485//p, 971) != 1 for p in (5, 97))
    assert pow(2, 48, 485) == 1
    assert all(pow(2, 48//p, 485) != 1 for p in (2, 3))
    f, ell = 1, 3
    for r in range(1, 197):
        if r >= 4 and (r-4) % 48 == 0:
            assert (f-16) % 971 == 0
        f, ell = f*ell % 971, (ell*ell-2) % 971
    f, ell = 3, 0
    for r in range(3, 30):
        f, ell = f*ell % 7, (ell*ell-2) % 7
        assert (f+4) % 7 == 4


if __name__ == "__main__":
    print("PASS: primitive lattice cases:", check_seeds())
    print("PASS: denominator divisibilities:", check_windows())
    check_lucas()
    print("PASS: Lucas valuations, seed 971, and nonreturning seed 7.")


\end{lstlisting}



\clearpage
\phantomsection
\addcontentsline{toc}{section}{References}
\begin{thebibliography}{99}


% Merge these entries into the root bibliography, omitting duplicate keys.
% All theorem/question/problem numbers used in the text refer to the
% explicitly cited arXiv versions, not necessarily to the journal numbering.

\bibitem{Saito2025}
K.~Saito, \emph{Mills' constant is irrational},
Mathematika \textbf{71} (2025), no.~3, e70027.
Questions~1.6 and~1.7 are numbered as in arXiv:2404.19461v2,
8 May 2025.
\url{https://arxiv.org/abs/2404.19461v2};
\url{https://doi.org/10.1112/mtk.70027}.

\bibitem{SaitoVariants}
K.~Saito, \emph{Transcendency of variants of Mills' constant},
The Ramanujan Journal \textbf{70} (2026), article~69,
published 10 August 2026.
\url{https://doi.org/10.1007/s11139-026-01443-0}.
Theorem numbers used here refer to arXiv:2508.16068v3,
7 December 2025, especially Theorems~1.7 and~1.8.
\url{https://arxiv.org/html/2508.16068v3}.

\bibitem{SaitoIntervals}
K.~Saito, \emph{Intervals without primes near an iterated
linear recurrence sequence}, International Journal of
Number Theory \textbf{22} (2026), no.~6, 1269--1279.
\url{https://doi.org/10.1142/S1793042126500685}.
Problem numbers~1.1, 1.7 and~1.8 refer to
arXiv:2504.14968v1, 21 April 2025.
\url{https://arxiv.org/html/2504.14968v1}.

\bibitem{Schneider1934}
T.~Schneider, \emph{Transzendenzuntersuchungen periodischer
Funktionen I. Transzendenz von Potenzen}, Journal f\"ur die
reine und angewandte Mathematik \textbf{172} (1934), 65--69.
\url{https://eudml.org/doc/149900}.

\bibitem{VMRC007}
Keio, \emph{Dyadic Lucas Prime Cascades and Composite
Fibonacci Shifts}, user-supplied research corpus,
manuscript [007], dated 5 September 2026.
The full all-shift Lucas/Fibonacci proof was already present
in this manuscript before the present reconstruction.

\bibitem{VMRC008}
Keio, \emph{Additive Clocks, Non-Reversible Recurrences, and
Transcendental Prime-Representing Constants}, user-supplied
research corpus, manuscript [008], dated 5 September 2026.
Source of the residue-surjectivity and non-reversible-clock
results reconstructed here.

\bibitem{VMRC010}
Keio, \emph{Arithmetic Constraints on a Hypothetically
Algebraic Mills Constant: Gap Divisibility, Transcendental
Exponents, and Cubic Reconstruction}, user-supplied research
corpus, manuscript [010], dated 5 September 2026.
Source of the conditional gap theorem.

\bibitem{VMRC042}
Keio, \emph{Cubic Pisot Floors: Composite Windows, Sparse
Exceptions, and Finite-Sieve Barriers}, user-supplied
research corpus, manuscript [042], dated 21 October 2026.
Source of the simultaneous three-adic composite-window
criterion. These internal manuscripts are cited for
provenance, not as independently peer-reviewed publications.


% Insert into the main thebibliography; retain just one copy of existing keys.
\bibitem{FR2026}
S.~Fischler and T.~Rivoal,
\emph{The ring of values of E-functions is not a field},
arXiv:2610.16603, 2026.
Author PDF dated 14 October 2026.
\url{https://www.imo.universite-paris-saclay.fr/~fischler/Ecroix.pdf}.

\bibitem{FR2016}
S.~Fischler and T.~Rivoal,
\emph{Arithmetic theory of E-operators},
Journal de l'\'Ecole polytechnique---Math\'ematiques
\textbf{3} (2016), 31--65.
DOI: \href{https://doi.org/10.5802/jep.28}{10.5802/jep.28}.
\url{https://www.imo.universite-paris-saclay.fr/~fischler/ateo.pdf}.

\bibitem{Beukers2006}
F.~Beukers,
\emph{A refined version of the Siegel--Shidlovskii theorem},
Annals of Mathematics (2) \textbf{163} (2006), no.~1, 369--379.
DOI: \href{https://doi.org/10.4007/annals.2006.163.369}{10.4007/annals.2006.163.369}.
\url{https://annals.math.princeton.edu/2006/163-1/p08}.

\bibitem{Lepetit2021}
G.~Lepetit,
\emph{Le th\'eor\`eme d'Andr\'e--Chudnovsky--Katz ``au sens large''},
North-Western European Journal of Mathematics \textbf{7} (2021), 81--149.
\url{https://nwejm.univ-lille.fr/index.php/nwejm/article/view/23}.
\url{https://nwejm.univ-lille.fr/index.php/nwejm/article/download/23/20/40}.



\bibitem{NISTDLMF}
NIST Digital Library of Mathematical Functions, Sections 5.7, 6.2,
16.17, and 16.21. Gamma expansions, exponential integrals,
and the Meijer $G$ integral and differential equation.
\url{https://dlmf.nist.gov/5.7}; \url{https://dlmf.nist.gov/6.2};
\url{https://dlmf.nist.gov/16.17}; \url{https://dlmf.nist.gov/16.21}.
Accessed 5 November 2026.

\bibitem{CrKo2025}
T.~Crmari\'c and V.~Kova\v{c},
\emph{On the irrationality of certain super-polynomially decaying series},
Colloquium Mathematicum \textbf{179} (2025), 55--68.
\url{https://doi.org/10.4064/cm9628-5-2025};
\url{https://arxiv.org/html/2504.18712v1}.

\bibitem{Lambrinoudis270}
C.~Lambrinoudis,
\emph{Algebraic Independence in the Affine Case of Erd\H{o}s Problem 270},
manuscript and companion repository, version 1.1.0, 24 August 2026.
\url{https://github.com/clambro/erdos-270-transcendence}.
Cited for the documented history of the base case; the broad affine
claims are not inputs to the present proofs.

\bibitem{CMF2026}
R.~Kalisch, T.~Monderer, E.~Leibtag, S.~Weinbaum, M.~Shalyt, and
I.~Kaminer, \emph{An Infinite Family of Ap\'ery-like Recurrences for
Euler's Constant}, arXiv:2610.13919v1, 9 October 2026.
\url{https://arxiv.org/html/2610.13919v1}.
Targets: Conjecture~1 and the analytic-identity component of Section~9.2(1).

\bibitem{Integrated2026}
Keio and Shaitetsu, \emph{Euler's Constant and Shifted Mills Constants:
Arithmetic Structure, Reciprocal Series, and Proof-Transfer Limits},
integrated working manuscript, 5 November 2026, 37 pages.
This is the preceding manuscript in the present research dialogue.

\bibitem{VMRC}
Keio, \emph{Verifiable Mathematical Reasoning Corpus},
working-manuscript repository, snapshot consulted 5 November 2026.
\url{https://github.com/QEEDDD/Verifiable-Mathematical-Reasoning-Corpus}.
The dated corpus entries above record provenance. They do not assert
independent peer review or establish historical priority.
\end{thebibliography}
\end{document}
